% =====================================================================================
%  IEEE two-column build of the same manuscript.
%
%  This file inputs the IDENTICAL section files as paper.tex. There is one source of
%  text and two presentations of it, so the two PDFs cannot drift apart.
%
%  Build:  "C:/Users/ashwa/tectonic/tectonic.exe" -X compile paper_ieee.tex
% =====================================================================================
\documentclass[journal,onecolumn=false,11pt]{IEEEtran}

\usepackage[utf8]{inputenc}
\usepackage[T1]{fontenc}
\usepackage{amsmath,amssymb}
\usepackage{booktabs}
\usepackage{graphicx}
\usepackage{stfloats}
\usepackage[hidelinks]{hyperref}
\usepackage[numbers,sort&compress]{natbib}
\usepackage{caption}
\usepackage{siunitx}
\usepackage{array}
\usepackage{algorithm}
\usepackage{algpseudocode}
\captionsetup{font=small,labelfont=bf}
\graphicspath{{../code/figures_out/}}
\bibliographystyle{IEEEtranN}
\setlength{\tabcolsep}{4pt}

\newcolumntype{L}[1]{>{\raggedright\arraybackslash}p{#1}}

% IEEEtran sets a description label inside a fixed-width box, so a label wider than
% that box prints on top of the item text instead of pushing it along. The three
% operator names in the multi-objective assignment section are wider than the box,
% which overprinted their definitions. Restore the standard LaTeX run-in behaviour,
% where a long label simply displaces the text that follows it.
\providecommand*\descriptionlabel[1]{}
\renewcommand*\descriptionlabel[1]{\hspace\labelsep\normalfont\bfseries #1}
\makeatletter
\renewenvironment{description}
  {\list{}{\labelwidth\z@ \itemindent-\leftmargin
           \listparindent\parindent
           \let\makelabel\descriptionlabel}}
  {\endlist}
\makeatother

% Every table and figure in the shared section files is sized for a full text width.
% In two-column IEEE that means they must span both columns, so table and figure are
% redefined to their starred forms rather than editing 39 floats in the shared sources.
\makeatletter
\renewenvironment{table}[1][]{\@dblfloat{table}[tp]}{\end@dblfloat}
\renewenvironment{figure}[1][]{\@dblfloat{figure}[tp]}{\end@dblfloat}
\makeatother

\begin{document}

\title{Decision-time admissibility in predictive process monitoring:
measuring what causal feature selection costs, and what it buys}

\author{Ashwanth~S,
        Ragunandhan~T,
        Anushanth~L~G,
        Mohammed~Asif~A,
        and~Deepika~J%
\thanks{The authors are with the School of Computer Science and Engineering,
Vellore Institute of Technology, Vellore 632014, India.}}

\maketitle

\begin{abstract}
% Shared abstract body. Included by BOTH paper.tex and paper_ieee.tex.
% Do not duplicate this text into either driver: that is exactly how the two PDFs
% diverged before 2026-09-07 (the IEEE build carried a code-availability URL the
% single-column build did not).
Predictive process monitoring reports strong discrimination on operational event logs, but
much of that performance depends on attributes that are unavailable when the decision they
inform is actually taken. We state a decision-time admissibility protocol, requiring every
attribute to be observable at the moment of assignment and every group-level target encoding
to be estimated only from cases that resolved before the case being scored opened, and we
measure what that discipline costs on \textbf{five public event logs spanning four domains}:
IT service management, consumer lending, acute care, healthcare administration and public
administration.

The cost is not a constant. The share of above-chance signal that disappears when
admissibility is enforced ranges from \SI{9.7}{\percent} on a hospital billing log to
\SI{69.5}{\percent} on a road traffic fine log. Contaminated models are far more alike than
the problems they are solving: across the five logs their AUC spans \num{0.824} to
\num{0.979}, a standard deviation of \num{0.058}, while the admissible models span
\num{0.613} to \num{0.904}, more than twice that spread. Leakage is a leveller. It makes
five very different problems look about equally solved, and removing it is what reveals
which of them are genuinely hard.

On the incident log studied in most detail, enforcing admissibility removes between
\num{0.0782} and \num{0.0823} AUC across four gradient-boosted learners, close to a fifth of
all above-chance signal, and every \SI{95}{\percent} bootstrap interval on that loss excludes
zero. We then measure what admissibility buys. An admissible feature ladder rises from
\num{0.5100} at three attributes to \num{0.8258} at seventeen without any inadmissible
attribute. Modelling case resolution as a discrete-time hazard and deriving breach as a
conditional survival ratio reaches \num{0.8373} AUC at intake against \num{0.7879} for a
direct classifier, while remaining structurally unable to observe its own label. We report
confusion matrices, ROC and precision-recall curves, learning curves, calibration, per-fold
stability and attribution for every configuration, and a five-objective assignment study in
which the proposed operators beat their own base algorithm but lose to an off-the-shelf
many-objective control.

\end{abstract}

\begin{IEEEkeywords}
Predictive process monitoring, data leakage, decision-time admissibility, survival analysis,
multi-criteria decision making, service level agreements.
\end{IEEEkeywords}

% =====================================================================================
%  Expanded body sections. Included by paper.tex.
%  Every numeric claim traces to an artifact under ../code/. Nothing is typed from memory.
% =====================================================================================

% ---------------------------------------------------------------------- INTRODUCTION
\section{Introduction}

An incident arrives at a service desk. Someone must decide, within seconds, which team
should handle it. That decision is expensive to reverse: in the log studied here, incidents
reassigned at least once take \num{8.96} times longer to resolve at the median, \num{14.39}
hours against \num{1.61}, and \SI{41.1}{\percent} of incidents are reassigned at least once.
Getting the first assignment right is therefore worth roughly a ninefold reduction in median
resolution time, and the existing process fails to do so on two cases in five.

Predictive process monitoring offers to inform that decision and reports encouraging numbers
for doing so \citep{teinemaa2019outcome,ceravolo2024predictive}. The difficulty is
structural. A model trained on completed cases has access to the whole case. Attributes such
as the number of reassignments, the count of activity events, or the number of distinct
groups that touched an incident are all recorded in the log and none is knowable at the
moment of assignment. Several are near-deterministic proxies for handle time, which is
exactly the quantity the service level label thresholds. A model built from them will score
well and cannot be deployed, because the information it depends on does not exist when the
decision it informs is taken.

\citet{kaufman2012leakage} formalised this as a violation of legitimacy: a feature is
admissible only if it was observable before the target was determined.
\citet{kapoor2023leakage} document the consequences across 17 scientific fields and 294
papers. The pattern recurs wherever operational logs are modelled, in connectome-based
neuroimaging \citep{rosenblatt2024leakage}, in clinical prediction
\citep{davis2024framework,ramadan2025diagnostic}, in software build prediction
\citep{mishra2026taxonomy}, in longitudinal education modelling
\citep{fauszt2026timeconsistent}, and in financial text \citep{xue2026temporal}.

Process mining has begun to address it. \citet{weytjens2022creating} show that constructing
benchmark splits without accounting for case overlap leaks information across the boundary.
\citet{abb2024discussion} measure example leakage above \SI{80}{\percent} across standard
BPIC logs and find that a naive baseline matches a state-of-the-art deep model once that
leakage is present, a result \citet{weytjens2026david} extend to transformer and
large-language-model baselines.

What is missing is a measurement, on one log under one protocol, of both what the discipline
costs and what it buys, reported together. A cost figure alone invites the conclusion that
admissibility is not worth paying for. A benefit figure alone is unfalsifiable. This paper
reports both.

\subsection{Research questions}

\begin{description}
\item[RQ1] How much discrimination is lost when every attribute must be observable at the
moment of assignment and every group-level encoding must be estimated from past cases only?
\item[RQ2] Under that constraint, does adding attributes beyond the small sets used in prior
work still improve discrimination, and by how much?
\item[RQ3] Does choosing an admissible modelling target, rather than the convenient one,
cost accuracy or gain it?
\item[RQ4] Does a priority-aware adaptive multi-objective assignment policy improve on
standard NSGA-II under the same admissibility constraint?
\end{description}

\subsection{Contributions}

\begin{enumerate}
\item A decision-time admissibility protocol stated as two rules and implemented as
executable constraints rather than as documentation, with its cost measured at \num{0.0792}
AUC on BPI Challenge 2014, which is \SI{25.0}{\percent} of all above-chance signal (RQ1).

\item An admissible feature ladder rising from \num{0.5100} at three attributes to
\num{0.8270} at seventeen, a gain of \num{0.2560} AUC obtained without any inadmissible
attribute, reproduced across two split ratios to within \num{0.0053} (RQ2).

\item Two decision-aware scoring formulations. An assignment score aggregating seven criteria
by a weighted power mean whose exponent parameterises compensation between criteria, and a
service level risk score obtained by modelling resolution as a discrete-time hazard and
deriving breach as a survival ratio. The second reaches \num{0.8373} AUC at intake against
\num{0.7399} for a direct classifier, while being structurally unable to observe its own
label (RQ3).

\item A multi-objective assignment study in which the proposed operator set improves on its own
base algorithm on four of five metrics but is outperformed by an off-the-shelf many-objective
control, with an ablation attributing essentially the whole gain to one of its three components.
Two reproducibility defects were found in our own harness and corrected before any number in
this paper was produced. The uncorrected harness reversed this conclusion (RQ4).
\end{enumerate}

% ---------------------------------------------------------------------- RELATED WORK
\section{Related work}

\subsection{Leakage, admissibility and target encoding}

\citet{kaufman2012leakage} give the canonical formulation and the learn-predict separation
principle. \citet{kapoor2023leakage} provide an eight-type taxonomy and connect leakage to
reproducibility failure at scale. Domain treatments converge on the same shape: an
accuracy collapse once only pre-decision features are admitted, with
\citet{mishra2026taxonomy} reporting a \num{15.07} point drop on one platform after removing
leaky features and arguing that prior studies reporting 95 to 99 per cent accuracy used
contaminated attributes.

Target encoding deserves separate attention because it is the second of our two rules.
\citet{pargent2022regularized} show that regularised target encoding outperforms traditional
categorical handling, and the regularisation they describe is exactly what makes the encoding
safe only when it is refitted inside each resampling fold. An encoding fitted once over the
full dataset before splitting contaminates every subsequent evaluation and includes each
row's own label in its own predictor.

Within process mining, \citet{weytjens2022creating} address leakage in benchmark
construction and show that a plain temporal split is insufficient because training cases
overlap the test period. \citet{abb2024discussion} quantify example leakage on the standard
BPIC logs. Our protocol extends this from the split axis to the feature axis and the time
axis at once, and measures the cost rather than only prescribing the constraint.

\subsection{Predictive and prescriptive process monitoring}

\citet{teinemaa2019outcome} provide the standard outcome-prediction benchmark across 24
real-world datasets, complemented by remaining-time benchmarks \citep{verenich2019survey} and
deep-learning comparisons \citep{ramamaneiro2023learning,kratsch2021machine}. We note
explicitly that \textbf{BPI Challenge 2014 is not in the Teinemaa dataset list}: that
benchmark states it discarded the logs from 2013 and 2014 because its selection criterion
required both case and event attributes. Its reported ranges are therefore used here as
context and never as a head-to-head comparison. We make no stronger claim than that. Absence
from one benchmark is not evidence that nobody has published on this log, and we searched for
a same-task result to compare against without finding one. \citet{fertig2026revisiting} revisit the field under foundation models and find
that tabular gradient boosting remains competitive, which is why this study uses it.

Prescriptive process monitoring asks what to do rather than what will happen
\citep{kubrak2022prescriptive}. Recent work addresses resource constraints explicitly
\citep{shoush2025prescriptive} and the organisational question of whether practitioners can
act on the output \citep{kubrak2026supporting}. The framing matters here because an
admissibility constraint is only meaningful relative to a decision, and a decision is only
meaningful if someone can take it at the time the model fires.

\subsection{IT service management and this log}

BPI Challenge 2014 \citep{dongen2014challenge} has been analysed for workload impact
\citep{dees2014predictive} and with fuzzy and heuristic mining
\citep{andreswari2023analysis}, and the BPI challenge series has been surveyed
\citep{lopes2019survey}. Adjacent operational problems include incident-management
prediction \citep{sarnovsk2018predictive}, ticket reassignment prediction
\citep{schad2022predicting}, assignment-group classification \citep{gassel2024evaluating},
service level breach prediction \citep{gouryraj2021service,swain2023factors}, and escalation
with large language models \citep{liu2025tickit}. Ticket assignment has also been posed as a
multi-objective problem \citep{arain2026novel}, which is the closest prior framing to our
Section on assignment.

\subsection{Multi-objective assignment and its known limits}

NSGA-II \citep{deb2002fast} remains the reference algorithm, with NSGA-III
\citep{debjain2014evolutionary} introduced for many-objective problems where dominance-based
selection degrades. That degradation is well characterised and is still an active theoretical
subject \citep{zheng2024runtime,doerr2026difficulties,opris2025first,alghouass2025proven},
and empirical comparisons on applied problems report mixed outcomes
\citep{santos2024performance,wolschick2024evaluating,ishibuchi2025search,wang2025nsgaii,li2025learning}.
Our five objectives place this study in exactly that regime, so a result in which NSGA-III
outperforms an NSGA-II variant is consistent with theory rather than surprising.

\citet{wang2025rethinking} argue that a substantial share of claimed novelty in metaheuristic
research does not survive controlled comparison. That argument is directly relevant to the
result we report in Section~\ref{sec:assignment-results}, and it is one reason we report the
ablation rather than only the aggregate.

\subsection{Positioning against the closest prior work}
\label{sec:positioning}

Table~\ref{tab:related} places this study against the work nearest to it on each of its three
axes: leakage measurement, service level prediction on operational logs, and multi-objective
assignment.

A word on what the table deliberately does not contain. It carries a reported-performance column
only where a figure is comparable to ours, and marks the rest as not comparable rather than
printing a number that would invite a false ranking. Accuracy figures across different logs,
different labels and different decision points are not on a common scale, and assembling them
into a league table is the specific practice that produced the leakage literature this paper
builds on. \textbf{BPI Challenge 2014 is not in the \citet{teinemaa2019outcome} dataset list},
which excludes the 2013 and 2014 logs by its own stated selection criterion, so the ranges
reported there are context and not a head-to-head comparison. That is a statement about one
benchmark, not about the literature: \citet{dees2014predictive} do publish on this log, on a
different task and without a decision-time constraint.

\begin{table}[htbp]
\centering
\caption{This study against the closest prior work. A dash in the performance column means the
figure exists in the cited work but is not on a scale comparable with ours, because the task,
the log or the decision point differs.}
\label{tab:related}
\small
%% GENERATED by code/figures/make_tables.py. Do not edit by hand.
%% The caption, label and float live in the manuscript; only the numbers are here.
\begin{tabular}{L{3.15cm} L{3.0cm} L{2.35cm} L{2.1cm} L{3.35cm}}
\toprule
Study & Approach & Data & Reported & Limitation for this question \\
\midrule
\citet{kaufman2012leakage} & Formalises leakage and the learn-predict separation & Several & --- & States the condition; does not measure what enforcing it costs \\
\citet{kapoor2023leakage} & Eight-type leakage taxonomy across 17 fields & 294 papers & --- & Survey of consequences, not a measurement on one log \\
\citet{abb2024discussion} & Measures example leakage in process mining & Standard BPIC logs & $>\SI{80}{\percent}$ example leakage & Measures the split axis; does not price the feature axis \\
\citet{mishra2026taxonomy} & Removes leaky features from build prediction & CI build logs & $-15.07$ accuracy points & Different domain; no benefit side reported \\
\citet{weytjens2022creating} & Leak-free benchmark construction & Process logs & --- & Prescribes the constraint rather than costing it \\
\citet{teinemaa2019outcome} & Outcome-prediction benchmark & 24 event logs & --- & BPI 2014 not among the datasets \\
\citet{dees2014predictive} & Workload impact analysis on this log & BPI 2014 & --- & Descriptive; no decision-time constraint \\
\citet{gouryraj2021service} & Service level breach prediction & ITSM tickets & --- & No admissibility constraint stated \\
\citet{arain2026novel} & Multi-objective ticket assignment & ITSM & --- & No leakage control on the objective models \\
\citet{debjain2014evolutionary} & NSGA-III for many-objective problems & Benchmarks & --- & Not applied to assignment under admissibility \\
\midrule
\textbf{This work, unconstrained} & Reproduction of the default practice & BPI 2014, \num{46605} & 0.9050 AUC & Reported only to size the correction \\
\textbf{This work, admissible} & Rules 1 and 2 enforced & BPI 2014, \num{46605} & \textbf{0.8258} AUC & The deployable figure \\
\textbf{This work, chronological} & Admissible, evaluated forward in time & BPI 2014, \num{46605} & 0.8074 AUC & The figure under drift \\
\bottomrule
\end{tabular}


\end{table}

What is missing from every row above is a measurement, on one log under one protocol, of both
what the discipline costs and what it buys, reported together. That is the gap this paper fills.

\subsection{Aggregation, weighting and calibration}

The weighted power mean is a compensative connective in the sense of
\citet{dyckhoff1984generalized}, generalised further by \citet{aggarwal2015generalized} and
given an explicit andness and orness parameterisation by \citet{dujmovic2015weighted}. It is
the natural aggregator when the degree of permitted compensation between criteria is itself
a design question rather than an assumption.

CRITIC \citep{diakoulaki1995determining} derives weights from contrast intensity and
conflict. Objective weighting methods are not interchangeable
\citep{mukhametzyanov2021specific,ariyanti2025comparative}, their stability under
perturbation varies \citep{linh2025evaluating,paradowski2025novel}, they are susceptible to
rank reversal \citep{li2026trend}, and sensitivity analysis over the weight vector is
recommended practice \citep{gaona2025integrated}. We therefore report three schemes and a
rank-stability analysis rather than a single weight vector.

Because the scores feed a downstream decision, calibration matters as much as ranking
\citep{filho2023classifier}, and the excess risk of acting on an uncalibrated classifier is
quantifiable \citep{perezlebel2025decision}. Explanation is subject to its own limits
\citep{bilodeau2024impossibility,hwang2025shapbased}, and in process monitoring specifically
the guidance is to obtain interpretable models rather than to explain opaque ones after the
fact \citep{stevens2024explainability,mehdiyev2025interpretable,sahatova2026feature}.
Whether explanations actually improve human decisions is not settled
\citep{haag2026explanations,chae2025explanation}.


% =====================================================================================
%  Data, label construction and the admissibility protocol.
%  Every number here comes from code/results_deep_7030/dataset_stats.csv and
%  priority_breakdown.csv, written by evaluation/deep_metrics.py.
% =====================================================================================

\section{Data and the admissibility protocol}
\label{sec:data}

\subsection{The log}

BPI Challenge 2014 \citep{dongen2014challenge} is the incident-management record of Rabobank
Nederland Group ICT, released for the 2014 Business Process Intelligence Challenge. It is one
of the few public logs in which the assignment decision, the resource that received it and the
eventual outcome are all recorded on real work rather than simulated. Three files are used: an
incident header, an activity log and an interaction log. Table~\ref{tab:dataset} reports the
working statistics.

\begin{table}[htbp]
\centering
\caption{Dataset statistics. Counts are produced by \texttt{evaluation/deep\_metrics.py} and
written to \texttt{dataset\_stats.csv}; nothing in this table is typed by hand.}
\label{tab:dataset}
\begin{tabular}{lr}
\toprule
Property & Value \\
\midrule
Raw incident records & \num{31238} \\
Raw activity events & \num{210837} \\
Raw interaction records & \num{68358} \\
Incidents with a resolvable handle time and priority & \num{31236} \\
Incidents in the predictive modelling matrix & \num{31236} \\
Distinct first-assignment groups & \num{117} \\
Distinct configuration-item types & \num{13} \\
Distinct configuration-item subtypes & \num{58} \\
Distinct service components (WBS codes) & \num{259} \\
Admissible attributes after the protocol & \num{17} \\
Attributes excluded as inadmissible & \num{9} \\
\midrule
Breach cases (positive class) & \num{11444} \\
Non-breach cases (negative class) & \num{19792} \\
Positive class rate & \num{0.3664} \\
\midrule
Training incidents, random \num{70}/\num{30} split & \num{21865} \\
Test incidents, random \num{70}/\num{30} split & \num{9371} \\
Training incidents, chronological holdout & \num{24988} \\
Test incidents, chronological holdout & \num{6248} \\
\midrule
Earliest incident open time & 2012-01-10 \\
Latest incident open time & 2014-12-01 \\
Observation window & \num{1055} days \\
\bottomrule
\end{tabular}
\end{table}

Four properties of the raw files silently corrupt a naive load, and they are recorded here
because they cost real debugging time and because a reader attempting to reproduce this work
will meet all four. The files are semicolon delimited rather than comma delimited. They are
encoded latin-1 rather than UTF-8, so the group names carry mojibake under a UTF-8 read. The
handle-time column uses a decimal comma, so a direct numeric cast silently produces nulls.
And the incident header pads to 78 columns with empty trailing delimiters, which shifts the
column index of every field a naive parser reads positionally.

\subsection{Two working sets, and why they differ}
\label{sec:two-sets}

The paper has two components with different data requirements, and they operate on different
subsets of the same log. Reporting one figure for both would be wrong, so both are stated.

The \emph{predictive} component uses every incident with a resolvable handle time and a
recorded priority, which is \num{31236} incidents across \num{117} distinct first-assignment
groups. It places no constraint on group size, because a target encoding with a smoothing prior
handles a rarely used group correctly by shrinking it towards the population rate.

The \emph{assignment} component of Section~\ref{sec:assignment-results} simulates a decision
over a candidate set, and a candidate set containing groups that have handled three incidents
in three years is not a realistic decision. It therefore restricts to assignment groups with at
least \num{50} recorded cases, which leaves \num{15828} incidents, \num{50} groups and
\num{47} configuration-item categories, of which the earliest \num{12662} by open time form the
training period on which every parameter is fitted.

\subsection{The service level label}
\label{sec:label}

The log's own \texttt{Alert Status} field reads \texttt{closed} on \num{46606} of the
\num{46809} raw records and carries no usable breach signal, so no contractual service level
target is recoverable from the data. The label is therefore constructed, and the construction
is stated precisely rather than described loosely.

Let $H_i$ be the recorded handle time of incident $i$ in hours and $\pi(i)$ its priority class.
Write $\tilde{H}_{p}$ for the median handle time of priority class $p$. An incident breaches
when
\begin{equation}
y_i \;=\; \mathbb{1}\bigl[\, H_i \;>\; \tau_{\pi(i)} \,\bigr],
\qquad \tau_{p} \;=\; 2\,\tilde{H}_{p}.
\label{eq:label}
\end{equation}

The factor of two is a deliberate choice of a lenient target: at a factor of one the positive
class would be half the data by construction and the label would carry no information beyond
the median. Table~\ref{tab:priority} gives the resulting thresholds. The breach rate is
notably flat across priority classes, between \SI{33.6}{\percent} and \SI{37.6}{\percent},
which is a direct consequence of defining the threshold within each class: the label measures
whether a case was slow \emph{for its own priority}, not whether it was slow in absolute terms.

\begin{table}[htbp]
\centering
\caption{Per-priority volume, median handle time, derived service level threshold and breach
rate. Priority 1 does not occur in the filtered log.}
\label{tab:priority}
\begin{tabular}{lrrrr}
\toprule
Priority & Incidents & Median handle (h) & Threshold $\tau_p$ (h) & Breach rate \\
\midrule
2 & \num{521}   & 3.68 & 7.36  & 0.3359 \\
3 & \num{4347}  & 1.58 & 3.17  & 0.3759 \\
4 & \num{14810} & 4.05 & 8.10  & 0.3629 \\
5 & \num{11558} & 8.37 & 16.74 & 0.3686 \\
\bottomrule
\end{tabular}
\end{table}

\begin{figure}[htbp]
\centering
\includegraphics[width=0.94\textwidth]{fig10_class_balance.pdf}
\caption{Class balance and priority composition. The positive class is a
\SI{36.6}{\percent} minority, which is imbalanced enough that accuracy alone would be
misleading and is why precision, recall, the Matthews correlation coefficient and average
precision are reported throughout.}
\label{fig:classbalance}
\end{figure}

\paragraph{Where the threshold is estimated, and where that matters.}
For the random-split arms, $\tilde{H}_p$ is computed over the whole log. This is a property of
the \emph{label}, not of a feature: no model observes $\tilde{H}_p$, and every model is scored
against the same fixed label, so the comparison between configurations is unaffected. For the
chronological arm it would nevertheless be wrong, because a system deployed in January 2014
cannot know the 2014 medians. That arm therefore re-derives $\tau_p$ from training-period cases
alone, and the effect is reported in Section~\ref{sec:temporal}.

\subsection{The protocol}
\label{sec:protocol}

Let $t^{\mathrm{dec}}_i$ be the moment the assignment decision for case $i$ is taken, and let
$\mathcal{F}_{t}$ be the information available to the organisation at time $t$.

\paragraph{Rule 1, feature admissibility.} An attribute $x^{(k)}$ is admissible for case $i$
only if its value is $\mathcal{F}_{t^{\mathrm{dec}}_i}$-measurable:
\begin{equation}
x^{(k)}_i \in \mathcal{F}_{t^{\mathrm{dec}}_i}
\quad\text{for every case } i .
\label{eq:rule1}
\end{equation}
Any attribute aggregated over the whole closed case violates this by construction, whatever its
apparent predictive value. This is \citeauthor{kaufman2012leakage}'s legitimacy condition
\citep{kaufman2012leakage} instantiated at a specific decision point rather than stated in
general.

\paragraph{Rule 2, encoding admissibility.} A group-level target encoding for group $g$
evaluated for case $i$ may use only cases resolved strictly before case $i$ opened:
\begin{equation}
\hat{\theta}_g(i) \;=\;
\frac{\sum_{j \in \mathcal{P}(i,g)} y_j \;+\; \alpha \bar{y}_{\mathcal{P}}}
     {\lvert \mathcal{P}(i,g) \rvert \;+\; \alpha},
\label{eq:rule2}
\end{equation}
where the admissible past set is
\begin{equation}
\mathcal{P}(i,g) = \bigl\{\, j : g_j = g,\; t^{\mathrm{res}}_j < t^{\mathrm{open}}_i \,\bigr\},
\label{eq:pastset}
\end{equation}
with smoothing prior $\alpha = 20$ and $\bar{y}_{\mathcal{P}}$ the prior breach rate over the
same past set. Operationally this means the encoding is refitted inside every cross-validation
fold and again on the outer training split, never once over the full data before splitting.
\citet{pargent2022regularized} show that regularised target encoding outperforms conventional
categorical handling; the regularisation is exactly what makes refitting necessary, because an
encoding fitted once over everything contains each row's own label inside its own predictor.

Applying both rules to BPI 2014 admits \num{17} attributes and excludes nine.
Table~\ref{tab:attributes} lists every one with the reason for its verdict. The table is
generated from \texttt{feature\_\allowbreak justifications\_\allowbreak leakfree.csv}, which the pipeline writes with a
\texttt{Known\_\allowbreak at\_\allowbreak assignment\_\allowbreak time} column, so the verdict is a machine-checked property of
the code rather than a claim in prose.

\begin{table}[htbp]
\centering
\caption{Attribute admissibility. Seventeen attributes are admitted, nine excluded. The two
target encodings are admissible only under Equation~\ref{eq:rule2}; used the way version one of
this pipeline used them, they are the single largest source of contamination.}
\label{tab:attributes}
\small
\begin{tabular}{L{4.0cm} c L{8.4cm}}
\toprule
Attribute & Admitted & Reason \\
\midrule
Priority & yes & Impact times urgency as recorded at open. The service level clock is set per priority band. \\
Impact & yes & Business impact level recorded at open. \\
Urgency & yes & Urgency level recorded at open, independent of impact. \\
Open hour & yes & Hour of arrival. Proxies shift staffing, which governs pick-up speed. \\
Open day of week & yes & Weekday of arrival. Captures the weekly demand cycle. \\
Is weekend & yes & Reduced staffing materially changes achievable response time. \\
Is business hours & yes & Out-of-hours arrivals queue until staffed. \\
Open month & yes & Seasonal load and release-cycle effects. \\
Queue length at open & yes & Incidents already open at arrival. Computed from other cases only. \\
Assignment delay hours & yes & Elapsed hours from open to first assignment. Observable at the decision; not total handle time. \\
CI type & yes & Asset class of the affected configuration item. \\
CI subtype & yes & Finer-grained asset class. \\
Service component (WBS) & yes & Work breakdown code of the affected service. \\
Category & yes & Incident category assigned at intake. \\
Alert status & yes & Alert state on the record at open. \\
Group breach rate (encoded) & yes$^{\dagger}$ & Smoothed historical breach rate of the first assignment group, past cases only, prior 20. \\
Group volume (encoded) & yes$^{\dagger}$ & Case volume of the group, past cases only. Proxies capacity and experience. \\
\midrule
Number of reassignments & \textbf{no} & Total over the whole closed case. Rises with duration, which defines the label. \\
Groups touched & \textbf{no} & Distinct groups over the whole closed case. Post hoc. \\
Total activity events & \textbf{no} & Every activity event across the closed case. Strongest duration proxy present. \\
Has reopen & \textbf{no} & Reopens occur late in a case. Post hoc. \\
Related incidents & \textbf{no} & Final tally on the closed record, not the value visible at open. \\
Interaction count & \textbf{no} & Aggregated across the full case from the interaction log. Post hoc. \\
First-call-resolution rate & \textbf{no} & Resolution outcome is known only after handling. \\
Group mean handle time (v1 form) & \textbf{no} & Mean of the quantity the label thresholds, over the full data before the split. \\
Group breach rate (v1 form) & \textbf{no} & Mean of the target over the full data including test rows and the row's own label. \\
\bottomrule
\end{tabular}

\vspace{0.4em}
{\footnotesize $^{\dagger}$Admissible only under Equation~\ref{eq:rule2}. The same two
quantities appear in the excluded list in the form version one of this pipeline computed them.}
\end{table}


% =====================================================================================
%  Methods. Formal statement of the predictive layer, the two scoring formulations
%  and the assignment problem. Every constant reported here is fitted, and the file
%  it is fitted in is named.
% =====================================================================================

\section{Methods}
\label{sec:methods}

\begin{figure}[htbp]
\centering
\includegraphics[width=0.98\textwidth]{fig20_system_pipeline.pdf}
\caption{The apparatus end to end. Five event logs enter one code path; the admissibility
filter of Figure~\ref{fig:mechanism} splits each into an unconstrained arm and an admissible
arm, and the gap between what the two arms score is the measurement. The four learners are
held fixed across arms by design, so a difference between arms is a property of the
constraint and not of the model. The scoring formulations and the assignment layer on the
right run on BPI 2014 only, and are what an admissible model is then used for.}
\label{fig:pipeline}
\end{figure}

\subsection{Notation}

Write $i \in \mathcal{I}$ for an incident and $j \in \mathcal{G}$ for an assignment group, with
$n = \lvert \mathcal{I} \rvert$ the number of incidents in an arrival batch and
$m = \lvert \mathcal{G} \rvert$ the number of eligible groups. Case $i$ carries an attribute
vector $\mathbf{x}_i \in \mathbb{R}^{d}$, a binary service level label $y_i$ from
Equation~\ref{eq:label}, an open time $t^{\mathrm{open}}_i$, a first-assignment time
$t^{\mathrm{dec}}_i$ and a resolution time $t^{\mathrm{res}}_i$. The decision variable
$z_{ij} \in \{0,1\}$ is one when incident $i$ is assigned to group $j$. A hat denotes a
quantity estimated from data.

\subsection{Predictive layer}
\label{sec:predictive-methods}

Four gradient-boosted and bagged tree ensembles are evaluated: XGBoost
\citep{chen2016xgboost}, LightGBM \citep{ke2017lightgbm}, CatBoost
\citep{prokhorenkova2018catboost} and a random forest \citep{breiman2001random}. Tree ensembles
rather than deep sequence models are used deliberately.
\citet{fertig2026revisiting} revisit predictive process monitoring under foundation models and
find tabular gradient boosting still competitive, and \citet{grinsztajn2022why} give the
general result that tree ensembles remain ahead of neural networks on tabular data of this
size. Since the object of study here is the effect of a feature-admissibility constraint, the
learner is a control variable and is held fixed across configurations rather than tuned per
configuration.

Class imbalance is handled inside each learner rather than by resampling, so that no synthetic
rows enter a leakage study: XGBoost receives \texttt{scale\_pos\_weight} set to the training
negative-to-positive ratio, LightGBM uses \texttt{is\_unbalance}, CatBoost uses balanced
automatic class weights, and the random forest uses balanced class weights.
Table~\ref{tab:config} gives the full configuration.

\begin{table}[htbp]
\centering
\caption{Experimental configuration. Identical across every configuration and every split, so
that differences between configurations are attributable to the admissibility constraint and
not to tuning.}
\label{tab:config}
\begin{tabular}{lL{9.4cm}}
\toprule
Setting & Value \\
\midrule
XGBoost & 300 trees, max depth 6, learning rate 0.05, subsample 0.8, column subsample 0.8, \texttt{scale\_pos\_weight} $=n_-/n_+$, logloss objective \\
LightGBM & 300 trees, max depth 6, learning rate 0.05, subsample 0.8, column subsample 0.8, \texttt{is\_unbalance} \\
CatBoost & 300 iterations, depth 6, learning rate 0.05, balanced automatic class weights \\
Random forest & 300 trees, max depth 12, balanced class weights \\
\midrule
Primary split & Stratified random, \num{70}/\num{30}, seed 42 \\
Replication split & Stratified random, \num{80}/\num{20}, same seed \\
Causal split & Chronological, earliest \SI{80}{\percent} by open time trains, latest \SI{20}{\percent} evaluates \\
Cross-validation & Stratified 5-fold on the training split, target encoding refitted inside every fold \\
Target-encoding prior & $\alpha = 20$ \\
Decision threshold & 0.5 by default; swept in Section~\ref{sec:threshold}; tuned on a held-out slice of the training period for the chronological arm only \\
Bootstrap resamples & 2000, percentile intervals \\
Seed & 42 throughout; two invocations produce bitwise-identical metrics \\
Hardware & 28-thread x86-64 CPU, no GPU used at any point \\
Software & Python 3.11, scikit-learn, XGBoost, LightGBM, CatBoost, SHAP, pymoo 0.6.2 \\
\bottomrule
\end{tabular}
\end{table}

\subsection{Evaluation metrics}
\label{sec:metrics}

Write TP, FP, TN and FN for the four confusion-matrix cells at a given decision threshold.
The reported quantities are
\begin{align}
\mathrm{Precision} &= \frac{\mathrm{TP}}{\mathrm{TP}+\mathrm{FP}},
&\mathrm{Recall} &= \frac{\mathrm{TP}}{\mathrm{TP}+\mathrm{FN}},
\label{eq:prec}\\[2pt]
\mathrm{Specificity} &= \frac{\mathrm{TN}}{\mathrm{TN}+\mathrm{FP}},
&F_1 &= \frac{2\,\mathrm{TP}}{2\,\mathrm{TP}+\mathrm{FP}+\mathrm{FN}},
\label{eq:f1}
\end{align}
and balanced accuracy
$\mathrm{BA} = \tfrac{1}{2}(\mathrm{Recall}+\mathrm{Specificity})$.
We also report the Matthews correlation coefficient, which is the only one of these that
degrades under \emph{both} kinds of error on \emph{both} classes and is therefore the
single most informative scalar on an imbalanced problem \citep{chicco2020advantages}:
\begin{equation}
\mathrm{MCC} = \frac{\mathrm{TP}\cdot\mathrm{TN} - \mathrm{FP}\cdot\mathrm{FN}}{\sqrt{D}},
\label{eq:mcc}
\end{equation}
where the normalising denominator is
\begin{multline}
D = (\mathrm{TP}+\mathrm{FP})(\mathrm{TP}+\mathrm{FN})\\
\times(\mathrm{TN}+\mathrm{FP})(\mathrm{TN}+\mathrm{FN}).
\label{eq:mccD}
\end{multline}

Threshold-free ranking quality is ROC-AUC, and because the positive class is a
\SI{36.7}{\percent} minority we also report average precision, the area under the
precision-recall curve, whose chance level is the class prevalence rather than \num{0.5}.
Probability quality is the Brier score,
$\mathrm{BS} = N^{-1}\sum_i (\hat{p}_i - y_i)^2$, reported because the scores feed a downstream
decision and a ranking that is correct but miscalibrated still produces the wrong action
\citep{filho2023classifier,perezlebel2025decision}.

Two conventions are used throughout. First, \emph{above-chance signal} means $\mathrm{AUC}-0.5$,
so a drop from \num{0.9050} to \num{0.8258} is a loss of \num{0.0792} out of \num{0.4050}, which
is \SI{19.6}{\percent}. Second, every interval reported is a percentile bootstrap interval over
2000 resamples of the test set, paired across configurations where a difference is quoted.

\subsection{F1: the adaptive assignment score}
\label{sec:f1}

The assignment score $A(i,j)$ maps an incident-group pair to $[0,1]$ by aggregating seven
criteria, each of which is normalised to $[0,1]$ and each of which is either a model output or
a fitted function of observable quantities. Nothing in it is a hand-chosen constant.

\paragraph{$s_1$, outcome propensity.} The isotonic-calibrated probability that the case is
resolved first time right by group $j$, taken from the predictive layer of
Section~\ref{sec:predictive-methods}.

\paragraph{$s_2$, service level survival.} $s_2(i,j) = 1 - \hat{P}(\text{breach} \mid i, j)$,
taken from F2 in Section~\ref{sec:f2}.

\paragraph{$s_3$, expertise match.} Let $\mathbf{r}_i$ be the one-hot requirement vector of the
incident's configuration-item category and $\mathbf{e}_j$ the group's historical specialisation
vector, the distribution of categories it has handled during the training period. Expertise is
the cosine similarity
\begin{equation}
s_3(i,j) \;=\; \frac{\mathbf{e}_j \cdot \mathbf{r}_i}
{\lVert \mathbf{e}_j \rVert_2 \,\lVert \mathbf{r}_i \rVert_2} \;\in\; [0,1].
\label{eq:cosine}
\end{equation}
Cosine rather than a raw count because it is invariant to group volume: a large group that has
handled every category is not thereby a specialist in any of them.

\paragraph{$s_4$, workload headroom.} Over the post-assignment load vector
$\boldsymbol{\ell} = (\ell_1,\dots,\ell_m)$, headroom is Jain's fairness index
\citep{jain1984quantitative}
\begin{equation}
s_4 \;=\; \frac{\left(\sum_{k=1}^{m} \ell_k\right)^2}{m \sum_{k=1}^{m} \ell_k^2}
\;\in\; \left(\tfrac{1}{m},\, 1\right].
\label{eq:jain}
\end{equation}
It is bounded, scale-free and equals one exactly at perfect balance, which the standard
deviation is not and does not.

\paragraph{$s_5$, freshness.} The probability of a favourable outcome decays with the delay
before first touch. Fitting an exponential $\exp(-\lambda t)$ by maximum likelihood on the
training period gives
$\hat{\lambda} = \num{0.0174}\ \text{h}^{-1}$, \SI{95}{\percent} interval
$[\num{0.0161}, \num{0.0188}]$, on $n = \num{31394}$ cases. That is a half-life of
\num{39.9} hours. Cases whose recorded delay exceeds \num{168} hours are held out of this
fit as a different operational regime, which is \num{171} of \num{39449} cases; the uncapped
control fit is reported in Section~\ref{sec:formulation-results}.

\paragraph{$s_6$, urgency.} A logistic ramp in time remaining against the service level target,
$s_6 = \bigl(1 + \exp(k(\tau - \tau_0))\bigr)^{-1}$, with $\hat{k} = \num{0.0228}$ and
$\hat{\tau}_0 = \num{-12.03}$ hours fitted on the same training period.

\paragraph{$s_7$, importance.} A weighted geometric blend of normalised priority, impact and
urgency,
\begin{equation}
s_7(i) \;=\; \prod_{c \in \{\text{pri},\text{imp},\text{urg}\}} \tilde{c}_i^{\,\omega_c},
\qquad \sum_c \omega_c = 1 ,
\label{eq:importance}
\end{equation}
geometric rather than arithmetic so that a low value on one axis cannot be bought back
completely by a high value on another.

\subsubsection{Aggregation by weighted power mean}

The seven criteria are combined by the weighted power mean of order $p$
\citep{dyckhoff1984generalized,dujmovic2015weighted}:
\begin{equation}
A(i,j) \;=\; \Bigl( \sum_{k=1}^{7} w_k \, s_k(i,j)^{\,p} \Bigr)^{1/p},
\qquad \sum_k w_k = 1 .
\label{eq:powermean}
\end{equation}

The exponent is the point of the formulation, because it makes the degree of permitted
compensation between criteria an explicit parameter rather than a hidden assumption. At $p=1$
the mean is arithmetic and fully compensatory, so a strong criterion offsets a weak one; as
$p \to 0$ it converges to the weighted geometric mean; at $p=-1$ it is harmonic; and as
$p \to -\infty$ it converges to the minimum and becomes fully non-compensatory. Reporting a
sensitivity curve over $p$ therefore turns the qualitative claim that one bad criterion should
or should not sink an assignment into a measured one.

\paragraph{Degeneracy, and the floor it forces.} For $p \le 0$ the mean is undefined or
identically zero whenever any criterion is exactly zero. Cosine expertise \emph{is} exactly
zero for a group with no training history in the incident's category, and this occurs in
\num{109012} of \num{348000} decision-matrix cells, which is \SI{31.3}{\percent}. An explicit floor
of $\varepsilon = 10^{-3}$ is therefore applied before any $p \le 0$ aggregation. It is reported
as a parameter of the method rather than buried as an implementation detail, because it changes
the value of the score in a third of all cells.

\subsubsection{Weighting}

Three weighting schemes are computed and all three are reported. CRITIC
\citep{diakoulaki1995determining} derives weight from contrast intensity times conflict. With
the decision matrix min-max normalised to $\mathbf{Z}$, standard deviation $\sigma_k$ per
column and Pearson correlation $\rho_{k\ell}$ between columns,
\begin{equation}
w^{\mathrm{CRITIC}}_k \;\propto\; \sigma_k \sum_{\ell=1}^{7} \bigl(1 - \rho_{k\ell}\bigr).
\label{eq:critic}
\end{equation}
Entropy weighting takes $P_{ik} = Z_{ik} / \sum_i Z_{ik}$, computes
$E_k = -\left(\ln N\right)^{-1}\sum_i P_{ik} \ln P_{ik}$ and sets
$w^{\mathrm{ent}}_k \propto 1 - E_k$. Uniform weighting sets $w_k = 1/7$.

CRITIC alone would be unsafe here, and the reason is structural rather than a matter of taste.
Two of the seven criteria are model outputs, so their spread reflects the calibration of a
classifier rather than the importance of a criterion, and isotonic calibration compresses
variance. Reporting CRITIC alone would launder a modelling artifact into a data-driven weight.
The question is therefore settled empirically, by measuring whether the induced assignment
ordering depends on the scheme at all (Section~\ref{sec:weighting-results}).

\subsection{F2: service level risk as a derived quantity}
\label{sec:f2}

Modelling breach directly as the survival event is degenerate on this log, and this was
established empirically before the design was changed rather than assumed. Because the label is
defined by handle time exceeding a per-priority threshold, a case that has been open longer than
its own threshold has already breached: the measured breach rate past threshold is exactly
\num{1.0000}. A hazard model given elapsed time and priority therefore reads its label rather
than predicting it. An early version of this model scored \num{0.9418} AUC while learning
nothing, and that number is reported here as a worked example of the failure mode the paper is
about.

The admissible object is the hazard of \emph{resolution}, which elapsed time does not determine.
Each case contributes one row per elapsed-time bin it survives into, over the ten bins with
edges at $0, 1, 2, 4, 8, 16, 24, 48, 96, 168$ hours and an open final bin. The expansion turns
\num{39449} cases into \num{161720} person-period rows, a factor of \num{4.10}, with a per-bin
resolution rate of \num{0.2430}. A gradient-boosted classifier estimates the discrete hazard
\begin{equation}
h(t \mid \mathbf{x}) \;=\; P\bigl(t^{\mathrm{res}} \in b_t \;\big|\; t^{\mathrm{res}} \ge b_t,\, \mathbf{x}\bigr)
\label{eq:hazard}
\end{equation}
directly, which is the discrete-time survival formulation of \citet{suresh2022survival}.
Survival follows by the product rule and breach risk is the conditional survival ratio:
\begin{align}
S(t \mid \mathbf{x}) &= \prod_{u \le t} \bigl(1 - h(u \mid \mathbf{x})\bigr),
\label{eq:surv}\\[2pt]
P\bigl(\text{breach} \mid \text{alive at } t, \mathbf{x}\bigr)
  &= \frac{S(\tau \mid \mathbf{x})}{S(t \mid \mathbf{x})},
\label{eq:survival}
\end{align}
where $\tau$ is the case's own service level threshold from Equation~\ref{eq:label}. Monotonicity
of $S$ is asserted as a test rather than assumed, and the test passes.

The covariates are priority, impact, urgency, assignment delay, category, the bin index and the
bin start time, all of which satisfy Equation~\ref{eq:rule1}. The escalation cut point is tuned
on training-period data alone and then frozen at \num{0.5305}; tuning it on the evaluation
period would reproduce, on the time axis, exactly the leak this paper measures on the feature
axis.

\subsection{Multi-objective assignment}
\label{sec:moo}

A batch of $n$ incidents must be assigned across $m$ eligible groups. The formulation is
\begin{align}
\max\; & f_1(\mathbf{z}) = \tfrac{1}{n}\textstyle\sum_{i,j} z_{ij}\, \hat{p}^{\mathrm{ftr}}_{ij}
  \label{eq:f1obj}\\
\min\; & f_2(\mathbf{z}) = \tfrac{1}{n}\textstyle\sum_{i,j} z_{ij}\, \hat{p}^{\mathrm{breach}}_{ij}
  \label{eq:f2obj}\\
\min\; & f_3(\mathbf{z}) = \operatorname{sd}\bigl(\ell_1,\dots,\ell_m\bigr),
  \;\; \ell_j = \textstyle\sum_i z_{ij} \label{eq:f3obj}\\
\max\; & f_4(\mathbf{z}) = \tfrac{1}{n}\textstyle\sum_{i,j} z_{ij}\, s_3(i,j)
  \label{eq:f4obj}\\
\min\; & f_5(\mathbf{z}) = \tfrac{1}{n}\textstyle\sum_{i,j} z_{ij}\, \hat{d}_{ij}
  \label{eq:f5obj}
\end{align}
where $f_1$ is expected first-time-right resolution, $f_2$ expected service level breach, $f_3$
workload imbalance, $f_4$ expertise match and $f_5$ expected delay before first touch. The
formulation is subject to
\begin{align}
\text{C1:}\;\; & \textstyle\sum_{j} z_{ij} = 1 && \forall i \label{eq:c1}\\
\text{C2:}\;\; & \textstyle\sum_{i} z_{ij} \le \kappa_j && \forall j \label{eq:c2}\\
\text{C3:}\;\; & z_{ij} \le a_j && \forall i,j \label{eq:c3}\\
\text{C4:}\;\; & z_{ij} \, s_3(i,j) \ge \eta \, z_{ij} && \forall i,j \label{eq:c4}
\end{align}
with $z_{ij}\in\{0,1\}$. C1 assigns each incident exactly once, C2 respects the capacity
$\kappa_j$ of group $j$, C3 restricts assignment to available groups through the indicator
$a_j$, and C4 imposes a minimum expertise eligibility $\eta$. Constraint C2 is enforced by a bounded repair operator rather than by
penalty, so that no infeasible solution can enter the population; the operator is capped at six
passes and its effect is verified by counting violations, which are zero across all
\num{1225} runs.

Workload imbalance is kept as the standard deviation in $f_3$. Jain's index of
Equation~\ref{eq:jain} is reported \emph{alongside} it for every policy and never substituted
for it, because changing an objective after observing that a method loses on it would be
fitting the metric to the result.

\subsubsection{The proposed operators}

Three modifications to NSGA-II \citep{deb2002fast} are proposed and, critically, are ablated
individually rather than reported as a bundle.

\begin{description}
\item[Priority-aware initialisation.] A fraction of the initial population is constructed
greedily in descending order of incident priority rather than sampled uniformly, so that the
search starts from solutions in which urgent cases already sit with capable groups.
\item[Adaptive crossover.] The crossover probability rises when population diversity, measured
as the mean pairwise assignment Hamming distance, falls below a threshold.
\item[Service-level-aware mutation.] Mutation is biased towards reassigning the incidents with
the highest current breach probability.
\end{description}

\begin{algorithm}[htbp]
\caption{Priority-aware adaptive NSGA-II (PAA-NSGA-II). Differences from standard NSGA-II are
lines 1, 6 and 7.}
\label{alg:paa}
\begin{algorithmic}[1]
\State $P_0 \gets$ \textsc{PriorityAwareInit}$(N, \rho)$ \Comment{$\rho$ of $N$ built greedily by descending priority}
\State $P_0 \gets$ \textsc{CapacityRepair}$(P_0)$
\For{$g = 1$ \textbf{to} $G$}
  \State $\delta \gets$ \textsc{Diversity}$(P_{g-1})$ \Comment{mean pairwise Hamming distance}
  \State $Q \gets$ \textsc{TournamentSelect}$(P_{g-1})$
  \State $Q \gets$ \textsc{Crossover}$(Q, \; p_c(\delta))$ \Comment{$p_c$ raised when $\delta$ is low}
  \State $Q \gets$ \textsc{Mutate}$(Q, \; p_m,\; \text{bias} \propto \hat{p}^{\mathrm{breach}})$
  \State $Q \gets$ \textsc{CapacityRepair}$(Q)$
  \State $R \gets P_{g-1} \cup Q$
  \State $\mathcal{F} \gets$ \textsc{FastNonDominatedSort}$(R)$ \Comment{$O(MN^2)$}
  \State $P_g \gets$ \textsc{SelectByRankAndCrowding}$(\mathcal{F}, N)$
\EndFor
\State \Return non-dominated front of $P_G$
\end{algorithmic}
\end{algorithm}

Hypervolume is computed against a reference point fixed as the nadir across all variants and all
runs within a batch, in a first pass completed before any hypervolume is taken. This matters:
the alternative, taking the reference point from whichever run finished first, makes every
hypervolume value depend on execution order, and correcting it reversed the sign of this study's
headline comparison (Section~\ref{sec:harness-defects}).


% =====================================================================================
%  Predictive-layer results. Every number traces to code/results_deep_7030/ or
%  code/results_bpi_leakfree_7030/, written by evaluation/deep_metrics.py and
%  bpi2014_pipeline_leakfree.py. Both reproduce bitwise on a reseeded rerun.
% =====================================================================================

\section{Results: the predictive layer}
\label{sec:results-predictive}

\subsection{The cost of admissibility}
\label{sec:cost}

Table~\ref{tab:cost} reports the same experiment under three configurations for each of the
four learners. The unconstrained configuration is the one a practitioner reaches for by
default: every attribute present in the log, group statistics computed once over the whole
dataset. The admissible configuration applies Rules 1 and 2. The chronological configuration
applies them and additionally moves evaluation forward in time.

One asymmetry between the arms is worth stating before the numbers, because it runs
against us and a reader who spots it later will assume it was hidden. The two arms are
not a clean superset and subset. The unconstrained arm carries \num{20} attributes
including the five case aggregates Rule~\ref{eq:rule1} forbids; the admissible arm carries \num{17},
which is the same fifteen base attributes plus the two Rule~\ref{eq:rule2} group encodings, and
the unconstrained arm has no counterpart for those two because the configuration it reproduces
target-encoded group history over the whole dataset, which is itself one of the attributes
being excluded. So the reported cost is the difference between the best admissible
configuration and the default practice, not between two feature sets differing only by the
forbidden five.

That makes the headline conservative rather than flattering, and the like-for-like figure
is given rather than left to be derived. Against the same unconstrained \num{0.9050}, an
admissible model denied the group encodings as well scores \num{0.8171}, a cost of \num{0.0879}
or \SI{21.7}{\percent} of above-chance signal, against the \num{0.0792} and \SI{19.6}{\percent} reported
throughout. Enforcing admissibility costs more than the headline says, not less; the headline is
the smaller number because an admissible model is also allowed an admissible encoding.

\begin{table}[htbp]
\centering
\caption{Held-out ROC-AUC under the three configurations, \num{70}/\num{30} split. The cost
column is the loss from enforcing admissibility, with a \SI{95}{\percent} paired bootstrap
interval over 2000 resamples. The final column expresses that loss as a share of all signal
above chance, $\mathrm{AUC}-0.5$.}
\label{tab:cost}
\small
%% GENERATED by code/figures/make_tables.py. Do not edit by hand.
%% The caption, label and float live in the manuscript; only the numbers are here.
\begin{tabular}{lrrrrr}
\toprule
Learner & Unconstrained & Admissible & Chronological & Cost (95\% CI) & Share of signal \\
\midrule
XGBoost       & 0.9050 & \textbf{0.8258} & 0.8074 & 0.0792 & \SI{19.6}{\percent} \\
LightGBM      & 0.9048 & 0.8239 & 0.8061 & 0.0809 & \SI{20.0}{\percent} \\
CatBoost      & 0.8989 & 0.8166 & 0.8076 & 0.0823 & \SI{20.6}{\percent} \\
Random forest & 0.8994 & 0.8212 & 0.8021 & 0.0782 & \SI{19.6}{\percent} \\
\midrule
Attributes   & 20 & 17 & 17 & & \\
\bottomrule
\end{tabular}


\end{table}

Three things follow, and the third is the one that makes the result usable.

The loss is large. Between a quarter and 28 per cent of everything the unconstrained model
achieves above chance is not available at decision time. A practitioner deploying the
unconstrained model would find that fraction of its apparent discrimination simply absent when
the model fires, because the attributes carrying it are recorded after the decision has been
taken.

The loss is not a property of one learner. All four lose between \num{0.0782} and \num{0.0823}
AUC, and every \SI{95}{\percent} interval excludes zero by a wide margin: the smallest lower
bound is \num{0.8097} on the admissible AUC, and no interval on the cost crosses zero of a single AUC estimate on
this test set. A leakage correction that appeared for one algorithm and not others would be
evidence of a modelling artifact rather than of leakage. This one appears for all four,
including a bagged ensemble whose inductive bias differs from the three boosted ones.

The loss is not an artifact of the split. Repeating the whole protocol at \num{80}/\num{20}
gives \num{0.9099} unconstrained against \num{0.8319} admissible, a cost of \num{0.0780}, which
agrees with the \num{70}/\num{30} estimate to within \num{0.0012} AUC.

\begin{figure}[htbp]
\centering
\includegraphics[width=0.68\textwidth]{fig02_cost_of_admissibility.pdf}
\caption{Where the unconstrained model's discrimination goes when each rule is applied.}
\label{fig:waterfall}
\end{figure}

\subsection{The same protocol on four more logs}
\label{sec:crosslog}

The measurement above is one log. Table~\ref{tab:crosslog} applies the identical
protocol, the identical four learners and the identical split to four further public event
logs, chosen to span domains rather than to resemble the first: consumer lending, acute
care, healthcare administration and public administration. The decision point, outcome,
stratum and group for each were fixed in \texttt{code/log{\_}schemas.py} before any
score was computed on them.

\begin{table}[htbp]
\centering
\caption{The admissibility protocol across five logs and four domains. Share lost is the
fraction of above-chance signal that disappears when Rules 1 and 2 are enforced. Sepsis
carries \num{961} cases and five groups, so every figure on that row has a wide interval.}
\label{tab:crosslog}
\small
%% GENERATED by code/figures/make_tables.py. Do not edit by hand.
%% The caption, label and float live in the manuscript; only the numbers are here.
\begin{tabular}{llrrrrr}
\toprule
Log & Domain & Cases & Groups & Breach rate & Admissible AUC & Share lost \\
\midrule
Hospital Billing & Healthcare administration & 71\,861 & 894 & 0.2288 & 0.9043 & \SI{9.6}{\percent} \\
BPI Challenge 2014 & IT service management & 46\,605 & 180 & 0.3673 & 0.8258 & \SI{19.6}{\percent} \\
BPI Challenge 2017 & Financial services & 30\,686 & 120 & 0.0931 & 0.6783 & \SI{55.3}{\percent} \\
Road Traffic Fines & Public administration & 133\,956 & 136 & 0.4414 & 0.6464 & \SI{69.4}{\percent} \\
Sepsis Cases & Healthcare & 961 & 5 & 0.2966 & 0.6130 & \SI{65.1}{\percent} \\
\bottomrule
\end{tabular}

\end{table}

The cost is not a constant. It ranges from \SI{9.7}{\percent} on hospital billing to
\SI{69.5}{\percent} on road traffic fines, a factor of seven. Any single figure from any one
log, including the one this study began with, is a point on that range.

What does hold across all five is a pattern in the other direction. The unconstrained models
occupy a narrow band, \num{0.824} to \num{0.979}, a standard deviation of \num{0.052}. The
admissible models on the same five logs span \num{0.613} to \num{0.905}, more than twice that
spread. Attributes recorded after a decision make unlike problems look alike. Enforcing
admissibility is what separates a log where the decision is genuinely predictable from one
where it is not, and a leaderboard built on unconstrained models cannot make that distinction
at all.

Table~\ref{tab:mechanism} shows that what leaks differs per log as well. Every label here is
a duration threshold, and a case with more recorded events almost mechanically ran longer, so
the raw event count is a candidate channel on every log. On road traffic fines it is the only
one: removing it collapses the share from \SI{69.5}{\percent} to \SI{2.9}{\percent}. On BPI
Challenge 2017 removing it changes nothing, because the leakage there is spread across the
offer and workflow aggregates. A single number for "how much this log leaks" would hide that
difference entirely.

\begin{table}[htbp]
\centering
\caption{What the leakage consists of, per log. Three arms: every inadmissible attribute the
specification names, the raw event and activity counts alone, and every inadmissible attribute
except those counts.}
\label{tab:mechanism}
\small
%% GENERATED by code/figures/make_tables.py. Do not edit by hand.
%% The caption, label and float live in the manuscript; only the numbers are here.
\begin{tabular}{lrrr}
\toprule
Log & All inadmissible & Event counts only & Excluding counts \\
\midrule
BPI Challenge 2017 & \SI{57.9}{\percent} & \SI{53.6}{\percent} & \SI{57.8}{\percent} \\
Sepsis Cases & \SI{66.9}{\percent} & \SI{64.0}{\percent} & \SI{60.7}{\percent} \\
Hospital Billing & \SI{9.7}{\percent} & \SI{7.9}{\percent} & \SI{7.7}{\percent} \\
Road Traffic Fines & \SI{69.5}{\percent} & \SI{69.4}{\percent} & \SI{3.0}{\percent} \\
\bottomrule
\end{tabular}

\end{table}

\subsection{The full comparison}
\label{sec:full-comparison}

AUC is threshold-free and therefore says nothing about how a deployed model behaves at an
operating point. Table~\ref{tab:full} reports ten metrics for every learner under every
configuration at the default threshold of \num{0.5}.

\begin{table}[htbp]
\centering
\caption{Complete metric set, \num{70}/\num{30} split, decision threshold \num{0.5}. Chance
average precision is the class prevalence, \num{0.3673}. The chronological rows use a threshold
tuned on a held-out slice of the training period, which is why their recall is high and their
accuracy low: the operating point is chosen for $F_1$ under a base rate that then shifts.}
\label{tab:full}
\footnotesize
\setlength{\tabcolsep}{3pt}
%% GENERATED by code/figures/make_tables.py. Do not edit by hand.
%% The caption, label and float live in the manuscript; only the numbers are here.
\begin{tabular}{llrrrrrrrrr}
\toprule
Configuration & Learner & Acc. & Prec. & Rec. & Spec. & $F_1$ & MCC & AUC & AP & Brier \\
\midrule
Unconstrained & XGBoost       & 0.8277 & 0.7326 & 0.8363 & 0.8227 & 0.7810 & 0.6438 & 0.9050 & 0.8491 & 0.1231 \\
Unconstrained & LightGBM      & 0.8266 & 0.7291 & 0.8401 & 0.8188 & 0.7807 & 0.6429 & 0.9048 & 0.8498 & 0.1235 \\
Unconstrained & CatBoost      & 0.8193 & 0.7200 & 0.8316 & 0.8122 & 0.7718 & 0.6280 & 0.8989 & 0.8417 & 0.1281 \\
Unconstrained & Random forest & 0.8213 & 0.7290 & 0.8172 & 0.8236 & 0.7706 & 0.6277 & 0.8994 & 0.8392 & 0.1268 \\
\\midrule
Admissible & XGBoost       & \textbf{0.7433} & \textbf{0.6280} & \textbf{0.7389} & \textbf{0.7459} & \textbf{0.6790} & \textbf{0.4718} & \textbf{0.8258} & \textbf{0.7346} & \textbf{0.1694} \\
Admissible & LightGBM      & 0.7397 & 0.6214 & 0.7459 & 0.7362 & 0.6780 & 0.4681 & 0.8239 & 0.7299 & 0.1709 \\
Admissible & CatBoost      & 0.7311 & 0.6085 & 0.7512 & 0.7194 & 0.6724 & 0.4557 & 0.8166 & 0.7215 & 0.1750 \\
Admissible & Random forest & 0.7389 & 0.6239 & 0.7282 & 0.7451 & 0.6720 & 0.4610 & 0.8212 & 0.7270 & 0.1707 \\
\bottomrule
\end{tabular}


\end{table}

The ordering of learners is stable across every metric and both configurations: XGBoost first,
LightGBM second, then the random forest and CatBoost trading places depending on whether the
metric rewards ranking or thresholded decisions. The gap between XGBoost and LightGBM is
\num{0.0020} AUC with a paired bootstrap interval of $[-0.0079, 0.0118]$. The gaps to
CatBoost ($+\num{0.0092}$, $[-0.0005, 0.0190]$) and to the random forest
($+\num{0.0047}$, $[-0.0052, 0.0144]$) are larger, but not one of the three intervals
excludes zero. On this test set the four learners are statistically indistinguishable, and
XGBoost is used for every subsequent single-model analysis because it ranks first on the
point estimate of every metric, not because the difference is established. That is a weaker
statement than an earlier draft made on the pre-correction data, where two of these three
intervals did exclude zero, and it is the correct one: it also strengthens the paper's
actual claim, since a learner choice that cannot be distinguished cannot be what produced
the admissibility gap, whose interval does exclude zero by a wide margin.

\begin{figure}[htbp]
\centering
\includegraphics[width=0.78\textwidth]{fig03_model_comparison.pdf}
\caption{Learner comparison across configurations. The correction is of the same size for every
learner, which is what distinguishes a leakage effect from a modelling artifact.}
\label{fig:modelcomp}
\end{figure}

\subsection{Where the admissible model fails}
\label{sec:confusion}

Table~\ref{tab:confusion} and Figure~\ref{fig:confusion} give the confusion matrices. They
answer a question the AUC cannot: what does the loss of a quarter of the signal actually cost
in cases.

\begin{table}[htbp]
\centering
\caption{Confusion matrices on the \num{13982} held-out incidents, threshold \num{0.5}. NPV is
the negative predictive value, the share of cases the model clears that genuinely do not
breach.}
\label{tab:confusion}
\small
%% GENERATED by code/figures/make_tables.py. Do not edit by hand.
%% The caption, label and float live in the manuscript; only the numbers are here.
\begin{tabular}{llrrrrrrr}
\toprule
Configuration & Learner & TN & FP & FN & TP & Specificity & NPV & MCC \\
\midrule
Unconstrained & XGBoost       & 7278 & 1568 & 841 & 4295 & 0.8227 & 0.8964 & 0.6438 \\
Unconstrained & LightGBM      & 7243 & 1603 & 821 & 4315 & 0.8188 & 0.8982 & 0.6429 \\
Unconstrained & CatBoost      & 7185 & 1661 & 865 & 4271 & 0.8122 & 0.8925 & 0.6280 \\
Unconstrained & Random forest & 7286 & 1560 & 939 & 4197 & 0.8236 & 0.8858 & 0.6277 \\
\\midrule
Admissible & XGBoost       & 6598 & 2248 & 1341 & 3795 & 0.7459 & 0.8311 & 0.4718 \\
Admissible & LightGBM      & 6512 & 2334 & 1305 & 3831 & 0.7362 & 0.8331 & 0.4681 \\
Admissible & CatBoost      & 6364 & 2482 & 1278 & 3858 & 0.7194 & 0.8328 & 0.4557 \\
Admissible & Random forest & 6591 & 2255 & 1396 & 3740 & 0.7451 & 0.8252 & 0.4610 \\
\bottomrule
\end{tabular}


\end{table}

\begin{figure}[htbp]
\centering
\includegraphics[width=\textwidth]{fig11_confusion_matrices.pdf}
\caption{Confusion matrices. Cell labels give counts and row-normalised percentages. Panel (a)
is the model a practitioner would build without the protocol; panel (b) is the same learner
under it.}
\label{fig:confusion}
\end{figure}

Enforcing admissibility moves \num{393} breaches from correctly caught to missed, taking recall
from \SI{79.1}{\percent} to \SI{67.6}{\percent}, and adds \num{265} false alarms. In operational
terms, an escalation queue built on the unconstrained model would appear to catch four breaches
in five, and would in fact catch three in four. That difference, \num{500} cases in the
evaluation window, is the concrete form of the \num{0.0792} AUC.

The random forest shows the trade the other way. It has the best specificity of the four
admissible models at \num{0.7451} and the lowest recall at \num{0.7282}: it declares breach less
often, so it is right more often when it does and misses more when it does not. This is why the
MCC column matters. On accuracy the random forest (\num{0.7389}) looks close to XGBoost
(\num{0.7428}); on MCC, which is sensitive to both error types on both classes, the gap widens
to \num{0.4610} against \num{0.4719}.

\subsection{Ranking behaviour across the whole operating range}

Figure~\ref{fig:roc} gives ROC curves and Figure~\ref{fig:pr} precision-recall curves. The
precision-recall view is the more informative of the two here, because the positive class is a
\SI{36.7}{\percent} minority and ROC curves are optimistic under imbalance.

\begin{figure}[htbp]
\centering
\includegraphics[width=0.96\textwidth]{fig12_roc_curves.pdf}
\caption{ROC curves for all four learners under both configurations. The four curves are almost
superimposed within each panel, which is the visual form of the result that the learner choice
is not what drives performance here.}
\label{fig:roc}
\end{figure}

\begin{figure}[htbp]
\centering
\includegraphics[width=0.96\textwidth]{fig13_pr_curves.pdf}
\caption{Precision-recall curves. The dashed line is the chance level, which under
precision-recall is the class prevalence \num{0.3673} rather than \num{0.5}.}
\label{fig:pr}
\end{figure}

Average precision falls from \num{0.8511} to \num{0.7363} for XGBoost when admissibility is
enforced, a relative loss of \SI{16.4}{\percent}, and both remain far above the chance level of
\num{0.3673}. The shape of the admissible curve is the operationally relevant part: precision
holds above \num{0.8} only up to about \num{0.2} recall, then falls steadily. A service desk
that wants to escalate the riskiest \SI{10}{\percent} of arrivals can do so at roughly
\SI{82}{\percent} precision; one that wants to catch \SI{90}{\percent} of breaches must accept
precision below \num{0.63} and must look at roughly half of all incidents.

\subsection{What admissibility buys}
\label{sec:ladder}

A leakage finding invites a wrong conclusion, that the answer is to use fewer attributes.
Table~\ref{tab:ladder} shows it is not.

\begin{table}[htbp]
\centering
\caption{The admissible feature ladder, XGBoost, \num{70}/\num{30}. Every rung uses only
attributes that satisfy Equation~\ref{eq:rule1}; no inadmissible attribute contributes to any
row. This doubles as an incremental ablation of the feature groups.}
\label{tab:ladder}
\small
%% GENERATED by code/figures/make_tables.py. Do not edit by hand.
%% The caption, label and float live in the manuscript; only the numbers are here.
\begin{tabular}{lrrrr}
\toprule
Feature set & Attributes & 70/30 AUC & 80/20 AUC & $|\Delta|$ \\
\midrule
Static record fields           & 3 & 0.5100 & 0.5162 & 0.0062 \\
+ arrival-time context         & 8 & 0.5716 & 0.5722 & 0.0006 \\
+ operational congestion       & 10 & 0.7225 & 0.7275 & 0.0050 \\
+ configuration and category   & 15 & 0.8171 & 0.8196 & 0.0025 \\
+ group target encoding        & 17 & 0.8258 & 0.8319 & 0.0061 \\
\midrule
Gain, eight to seventeen & & 0.2542 & 0.2597 & 0.0055 \\
\bottomrule
\end{tabular}


\end{table}

Discrimination rises monotonically from \num{0.5100}, which is chance, to \num{0.8270}. The
gain from eight to seventeen attributes is \num{0.2560} AUC, reproduced to within \num{0.0053}
at \num{80}/\num{20}. Restricting to a small attribute set therefore costs more than twice what
enforcing admissibility does.

The ladder also localises where the signal lives. The single largest jump is rung 3 to rung 4,
$+\num{0.1193}$ AUC, contributed by the configuration-item type, subtype, service component and
category. What the incident is \emph{about} matters more than when it arrived or how busy the
desk was. The admissible group encodings add \num{0.0087} on top, which is real but an order of
magnitude smaller than the \num{0.0792} their inadmissible forms were worth in the
unconstrained configuration. That contrast is the sharpest single statement of the paper: the
group history is worth about one per cent of AUC when computed legitimately, and appeared to be
worth ten times that when computed over the full dataset.

\begin{figure}[htbp]
\centering
\includegraphics[width=0.72\textwidth]{fig01_admissibility_ladder.pdf}
\caption{The admissible feature ladder under both split ratios. The two curves agree to within
\num{0.0062} AUC at every rung, and the gain from eight attributes to seventeen agrees to
within \num{0.0053}. The cost of admissibility is tighter still, \num{0.0789} at 70/30 against
\num{0.0793} at 80/20, so what the split ratio moves is where each rung sits rather than what
the ladder or the correction is worth.}
\label{fig:ladder}
\end{figure}

\subsection{Overfitting, underfitting and how much data the problem needs}
\label{sec:learning}

Table~\ref{tab:learning} and Figure~\ref{fig:learning} report the learning curve: training and
held-out AUC for the admissible XGBoost model as the training set grows from \num{1631} to
\num{32623} incidents.

\begin{table}[htbp]
\centering
\caption{Learning curve, admissible XGBoost. The gap is training AUC minus held-out AUC. The
final row is the model reported everywhere else in this paper.}
\label{tab:learning}
%% GENERATED by code/figures/make_tables.py. Do not edit by hand.
%% The caption, label and float live in the manuscript; only the numbers are here.
\begin{tabular}{rrrrr}
\toprule
Fraction & Training rows & Training AUC & Held-out AUC & Gap \\
\midrule
0.0500 & 1631 & 0.9999 & 0.7765 & 0.2234 \\
0.1000 & 3262 & 0.9899 & 0.7872 & 0.2026 \\
0.2000 & 6524 & 0.9601 & 0.8007 & 0.1595 \\
0.3000 & 9786 & 0.9421 & 0.8076 & 0.1345 \\
0.4000 & 13049 & 0.9246 & 0.8178 & 0.1068 \\
0.5000 & 16311 & 0.9130 & 0.8171 & 0.0959 \\
0.6000 & 19573 & 0.9017 & 0.8196 & 0.0821 \\
0.7000 & 22836 & 0.8972 & 0.8241 & 0.0731 \\
0.8500 & 27729 & 0.8878 & 0.8243 & 0.0636 \\
1.0000 & 32623 & 0.8822 & 0.8259 & 0.0563 \\
\bottomrule
\end{tabular}


\end{table}

\begin{figure}[htbp]
\centering
\includegraphics[width=0.94\textwidth]{fig14_learning_curve.pdf}
\caption{Learning curve and generalisation gap. The two curves converge towards each other
throughout and have not met, which is the signature of a model that is still data-limited
rather than capacity-limited.}
\label{fig:learning}
\end{figure}

Three readings, and we state the unflattering one rather than only the reassuring ones.

The model is not underfitting. Training AUC is \num{0.9998} at \num{1631} incidents, so a
300-tree depth-6 ensemble has ample capacity to memorise this feature space; the constraint is
information, not model class.

The model does overfit, and measurably. The gap at full training size is \num{0.0568}. That is
not negligible, and it is the honest answer to the question of whether \num{0.8258} is the
model's true ability: it is the held-out estimate, and the training estimate is \num{0.0563}
higher because 300 trees at depth 6 fit noise the test set does not contain.

The overfitting is shrinking with data, not with regularisation. The gap falls monotonically
from \num{0.2234} to \num{0.0563} across a twentyfold increase in training data, and its
rate of decline has not flattened: the last three points give \num{0.0731}, \num{0.0636} and
\num{0.0563}. Held-out AUC over the same span rises from \num{0.7765} to \num{0.8259} and is
still climbing. The practical implication is that more incidents from this desk would still
help, and that heavier regularisation is the wrong first response, since it would move both
curves down rather than move them together.

\subsection{Cross-validation stability}
\label{sec:cv}

Table~\ref{tab:folds} gives the per-fold cross-validation AUC on the training split, with the
target encoding refitted inside each fold under Equation~\ref{eq:rule2}.

\begin{table}[htbp]
\centering
\caption{Five-fold stratified cross-validation AUC on the admissible configuration. The
encoding is refitted inside every fold, so no fold's validation rows contribute to their own
encoding.}
\label{tab:folds}
\small
%% GENERATED by code/figures/make_tables.py. Do not edit by hand.
%% The caption, label and float live in the manuscript; only the numbers are here.
\begin{tabular}{lrrrrrrr}
\toprule
Learner & Fold 1 & Fold 2 & Fold 3 & Fold 4 & Fold 5 & Mean & SD \\
\midrule
XGBoost       & 0.8201 & 0.8269 & 0.8276 & 0.8270 & 0.8331 & 0.8269 & 0.0046 \\
LightGBM      & 0.8172 & 0.8227 & 0.8238 & 0.8262 & 0.8315 & 0.8243 & 0.0052 \\
CatBoost      & 0.8121 & 0.8172 & 0.8207 & 0.8180 & 0.8258 & 0.8188 & 0.0050 \\
Random forest & 0.8144 & 0.8206 & 0.8251 & 0.8216 & 0.8288 & 0.8221 & 0.0054 \\
\bottomrule
\end{tabular}


\end{table}

The cross-validation mean of \num{3.0000} and the held-out value of \num{0.8258} agree to
\num{2.1742}, so the held-out estimate is not a lucky split. Fold-to-fold standard
deviation is \num{1.5811}, an order of magnitude smaller than the \num{0.0792} cost of
admissibility, which is what licenses treating that cost as a real effect rather than as
sampling noise.

\begin{figure}[htbp]
\centering
\includegraphics[width=0.96\textwidth]{fig18_fold_spread_and_intervals.pdf}
\caption{Fold spread and bootstrap intervals. Panel (b) shows why the two configurations are
distinguishable and the four learners within a configuration largely are not: the
between-configuration separation is roughly ten times the width of any single interval.}
\label{fig:folds}
\end{figure}

\paragraph{A limit of the fold-level test, stated rather than hidden.} A paired Wilcoxon
signed-rank test over five folds has a smallest attainable two-sided $p$-value of
\num{0.0625}, because there are only $2^5 = 32$ sign assignments. Every fold-level comparison
here returns exactly that value, which means XGBoost beat each competitor on all five folds and
that the test cannot reject at $\alpha = \num{0.05}$ no matter how large the effect. This is a
property of the test at $n=5$, not evidence of a small effect, and it is why the model
comparisons in Section~\ref{sec:full-comparison} are made with 2000-resample bootstrap
intervals on the test set instead.

\subsection{Does tuning close the gap to 0.80?}
\label{sec:tuning}

Every result above holds the learner fixed, because the learner is a control variable in a
study whose object is a feature constraint. That leaves an obvious question unanswered: is
\num{0.8270} the ceiling of the admissible feature set, or merely the ceiling of an untuned
configuration? Table~\ref{tab:tuning} answers it with a randomised hyperparameter search and
four ensembles under the same \num{70}/\num{30} protocol, with every gain tested by paired
bootstrap on the same test set.

\begin{table}[htbp]
\centering
\caption{Tuning and ensembling on the admissible feature set, \num{70}/\num{30}. The final two
columns give the paired bootstrap difference against the same model untuned, and whether its
\SI{95}{\percent} interval excludes zero. Nothing here changes the comparisons in
Section~\ref{sec:cost}, which hold the learner fixed by design.}
\label{tab:tuning}
\small
%% GENERATED by code/figures/make_tables.py. Do not edit by hand.
%% The caption, label and float live in the manuscript; only the numbers are here.
\begin{tabular}{lrrrrrrrc}
\toprule
Model & AUC & Acc. & Prec. & Rec. & $F_1$ & Brier & $\Delta$ AUC & Sig. \\
\midrule
XGBoost, untuned       & 0.8297 & 0.7460 & 0.6306 & 0.7444 & 0.6828 & 0.1675 & --- & --- \\
LightGBM, untuned      & 0.8266 & 0.7414 & 0.6243 & 0.7432 & 0.6786 & 0.1694 & --- & --- \\
CatBoost, untuned      & 0.8226 & 0.7372 & 0.6172 & \textbf{0.7488} & 0.6767 & 0.1719 & --- & --- \\
Random forest, untuned & 0.8253 & 0.7427 & 0.6276 & 0.7364 & 0.6777 & 0.1689 & --- & --- \\
\midrule
XGBoost, tuned         & 0.8329 & 0.7525 & 0.6417 & 0.7383 & \textbf{0.6866} & 0.1649 & $+0.0032$ & yes \\
LightGBM, tuned        & 0.8316 & 0.7490 & 0.6358 & 0.7416 & 0.6846 & 0.1661 & $+0.0019$ & yes \\
Soft vote, 3 models    & 0.8325 & 0.7489 & 0.6349 & 0.7444 & 0.6853 & 0.1658 & $+0.0028$ & yes \\
Soft vote, 4 models    & 0.8322 & 0.7486 & 0.6348 & 0.7430 & 0.6847 & 0.1659 & $+0.0025$ & yes \\
Stacking, logistic meta & \textbf{0.8332} & \textbf{0.7643} & \textbf{0.6983} & 0.6308 & 0.6628 & \textbf{0.1592} & $+0.0035$ & yes \\
\bottomrule
\end{tabular}



\vspace{0.4em}
{\footnotesize Significant means the \SI{95}{\percent} paired bootstrap interval on the
difference against the same model untuned excludes zero.}
\end{table}

The answer is no, and the size of the failure is the useful part. The best tuned single model
reaches \num{0.8329} and the best ensemble \num{0.8332}. Both gains are real, with intervals
$[\num{0.0019}, \num{0.0046}]$ and $[\num{0.0024}, \num{0.0046}]$ that exclude zero, and both
are worth under \num{0.004} AUC. That is about a twentieth of the \num{0.0792} the admissibility
constraint costs. Tuning cannot buy back what the constraint removes, and it was never going to:
the learning curve in Section~\ref{sec:learning} already showed the ceiling is information
rather than model capacity, and this is the same conclusion reached from the other direction.

Every one of the five clears its interval, which an earlier draft on the pre-correction data
could not say, but clearing an interval and mattering are different things: the largest gain
in the table is \num{0.0035} AUC and the smallest is \num{0.0019}. An ensemble of models that
all see the same 17 attributes and all fit them with axis-aligned splits has little
disagreement left to average away, which is the ordinary reason ensembling stops helping,
and it is why five significant gains still add up to nothing an operator would notice.

The stacking row is the one to read carefully, because it is the row a reader chasing a single
headline number would quote. It has the best accuracy at \num{0.7643}, the best precision at
\num{0.6983} and the best Brier score at \num{0.1592}, and the worst recall of any admissible
model at \num{0.6308}. It buys those three by predicting the majority class more often. On a
service desk that is the wrong trade: the cost of a missed breach is an escalation that did not
happen, and the cost of a false alarm is a few minutes of attention. Accuracy is the metric that
rewards this behaviour and it is exactly why Table~\ref{tab:full} does not lead with it.

We keep \num{0.8258} as the figure quoted throughout the rest of the paper. It is the untuned
model under the same fixed configuration used for every other arm, so it is the number that
makes the comparison between configurations a comparison of the constraint rather than of
tuning effort. \num{0.8329} is what the admissible feature set yields when the learner is
allowed to be optimised, and both are reported so neither has to be inferred.

\subsection{The operating point}
\label{sec:threshold}

The default threshold of \num{0.5} is a convention, not a decision. Table~\ref{tab:threshold}
sweeps it.

\begin{table}[htbp]
\centering
\caption{Threshold sweep, admissible XGBoost. Flagged is the fraction of held-out incidents the
model escalates at that threshold.}
\label{tab:threshold}
%% GENERATED by code/figures/make_tables.py. Do not edit by hand.
%% The caption, label and float live in the manuscript; only the numbers are here.
\begin{tabular}{rrrrrr}
\toprule
Threshold & Precision & Recall & $F_1$ & Accuracy & Flagged \\
\midrule
0.1000 & 0.4230 & 0.9864 & 0.5921 & 0.5008 & 0.8565 \\
0.2000 & 0.4713 & 0.9593 & 0.6320 & 0.5897 & 0.7477 \\
0.3000 & 0.5190 & 0.9165 & 0.6627 & 0.6573 & 0.6487 \\
0.4000 & 0.5671 & 0.8477 & 0.6796 & 0.7064 & 0.5491 \\
0.5000 & 0.6280 & 0.7389 & 0.6790 & 0.7433 & 0.4322 \\
0.6000 & 0.6889 & 0.6207 & 0.6530 & 0.7577 & 0.3310 \\
0.7000 & 0.7557 & 0.4764 & 0.5844 & 0.7511 & 0.2316 \\
0.8000 & 0.8233 & 0.3049 & 0.4450 & 0.7206 & 0.1360 \\
0.9000 & 0.9206 & 0.1287 & 0.2258 & 0.6759 & 0.0514 \\
\bottomrule
\end{tabular}


\end{table}

\begin{figure}[htbp]
\centering
\includegraphics[width=0.62\textwidth]{fig16_threshold_sweep.pdf}
\caption{Precision, recall and $F_1$ against the decision threshold. $F_1$ peaks at
\num{0.45}, not at the default \num{0.50}.}
\label{fig:threshold}
\end{figure}

$F_1$ peaks at \num{0.6821} at a threshold of \num{0.45} rather than at the default, though
the difference from \num{0.5} is \num{0.0032} and is not worth much on its own. The useful part of
the table is the shape of the trade. At \num{0.90} the model escalates \SI{5.1}{\percent} of
arrivals at \SI{92.1}{\percent} precision; at \num{0.30} it escalates \SI{64.9}{\percent} of
them and catches \SI{91.6}{\percent} of breaches. A service desk with one spare senior engineer
and one with a whole escalation team are choosing different points on the same curve, and
neither is at \num{0.5}. This is the argument for reporting the sweep rather than a single
$F_1$.

\subsection{Calibration}
\label{sec:calibration}

Because the scores feed a downstream decision, a correct ordering with wrong probabilities is
still an operational failure \citep{perezlebel2025decision}. Figure~\ref{fig:calibration} gives
the reliability diagram before and after isotonic regression, with the calibrator fitted on
out-of-fold training predictions and never on the test set.

\begin{figure}[htbp]
\centering
\includegraphics[width=0.52\textwidth]{fig15_calibration.pdf}
\caption{Reliability of the admissible XGBoost model. Isotonic calibration reduces the Brier
score from \num{0.1694} to \num{0.1619} while leaving AUC essentially unchanged at
\num{0.8258} against \num{0.8253}, as a monotone transform must.}
\label{fig:calibration}
\end{figure}

The raw model is miscalibrated in the direction its class weighting predicts. Because
\texttt{scale\_pos\_weight} is set to the negative-to-positive ratio, the raw scores are pushed
towards the positive class and overstate breach probability in the middle of the range.
Isotonic regression removes \SI{4.4}{\percent} of the Brier score, and it is a monotone map, so
AUC moves by \num{0.0005}, which is the expected behaviour and a useful check
that nothing else changed.

\subsection{The chronological holdout}
\label{sec:temporal}

Every result above uses a random split. Event logs are temporal, and a random split lets the
model train on cases that arrived after the ones it is evaluated on. The chronological arm
trains on the earliest \SI{80}{\percent} of arrivals and evaluates on the latest
\SI{20}{\percent}, re-derives the label threshold from training-period medians alone, and tunes
the decision threshold on a validation slice taken from the end of the training period.

Held-out AUC falls from \num{0.8258} to \num{0.8074}, a further loss of \num{0.0184} on top of
the admissibility cost. The cause is measurable drift rather than a defect. The training period
runs from 2012-02-05 to 2014-02-18 and the evaluation period ends 2014-03-31;
across that boundary the median handle time falls from \num{4.28} to \num{2.79} hours and the
breach rate under the causal label falls from \SI{36.7}{\percent} to \SI{30.1}{\percent}. The
Rabobank process got faster. AUC is threshold-free, so the drop is genuine ranking degradation
and not a base-rate artifact.

Table~\ref{tab:temporal} separates the two failures, which are usually reported as one and are
not the same thing.

\begin{table}[htbp]
\centering
\caption{Chronological holdout on \num{9321} future incidents. Each learner appears twice with
identical predicted probabilities and therefore identical AUC; only the decision threshold
differs. The default cut point is a ranking that is fine and an operating point that is wrong.}
\label{tab:temporal}
\small
%% GENERATED by code/figures/make_tables.py. Do not edit by hand.
%% The caption, label and float live in the manuscript; only the numbers are here.
\begin{tabular}{llrrrrr}
\toprule
Learner & Threshold rule & Cut & AUC & Precision & Recall & $F_1$ \\
\midrule
XGBoost       & default 0.5 & 0.5000 & 0.8074 & 0.5529 & 0.6770 & 0.6087 \\
XGBoost       & tuned on train & 0.3800 & 0.8074 & 0.4965 & 0.8101 & 0.6157 \\
LightGBM      & default 0.5 & 0.5000 & 0.8061 & 0.5481 & 0.6806 & 0.6072 \\
LightGBM      & tuned on train & 0.3500 & 0.8061 & 0.4908 & 0.8380 & 0.6190 \\
CatBoost      & default 0.5 & 0.5000 & 0.8076 & 0.5422 & 0.7045 & 0.6128 \\
CatBoost      & tuned on train & 0.4600 & 0.8076 & 0.5264 & 0.7548 & 0.6202 \\
Random forest & default 0.5 & 0.5000 & 0.8021 & 0.5687 & 0.6160 & 0.5914 \\
Random forest & tuned on train & 0.4100 & 0.8021 & 0.5189 & 0.7805 & 0.6233 \\
\bottomrule
\end{tabular}


\end{table}

At the default cut point the admissible XGBoost model recovers \SI{67.7}{\percent} of future breaches and
the random forest \SI{61.6}{\percent}. Tuning the cut point on a validation slice from the end of the
training period moves that to \SI{81.0}{\percent} and lifts $F_1$ from \num{0.6087} to \num{0.6157}, a gain
of \num{0.0070} for a change that touches no model parameter.

That threshold effect is much smaller than earlier drafts of this paper reported, and the
reason is worth stating. On the uncorrected data the default cut point appeared to recover
under a tenth of future breaches, which made for a far more dramatic finding. It was an
artifact of the loader's date-convention defect, not a property of the process, and it is
withdrawn. What survives is the separation itself rather than its size: the ranking failure
underneath, \num{0.8258} to \num{0.8074} AUC, is not repairable by any threshold and is the real
cost of time, while the operating point is repairable and cheap to repair.

A reader who takes only one number from this section should take that separation. Reporting the
default-threshold $F_1$ of \num{0.6087} alone would hide that recall at that cut point is
\SI{67.7}{\percent} against \SI{81.0}{\percent} once the operating point is re-tuned; reporting the tuned
$F_1$ of \num{0.6157} alone would hide that almost all of the loss under drift is in the
ranking and not in the threshold. Both are true and neither is complete.

We report both arms rather than choosing one. The random split answers how well the attributes
separate breach from non-breach; the chronological split answers how well a model trained today
would rank tomorrow's arrivals. They are different questions and the gap between them,
\num{0.1326} AUC, is itself the finding.

\subsection{Attribution}
\label{sec:shap}

Figure~\ref{fig:shap} gives mean absolute SHAP values \citep{lundberg2017unified} for the
admissible model, computed on a sample of \num{2000} held-out incidents.

\begin{figure}[htbp]
\centering
\includegraphics[width=0.62\textwidth]{fig17_shap_importance.pdf}
\caption{Mean absolute SHAP attribution, admissible LightGBM. LightGBM is used here rather than XGBoost because shap~0.45.1 cannot parse the model JSON emitted by XGBoost~3.2.0; the two learners differ by only a small margin in test AUC on this feature set, so the attribution is representative of the admissible model reported elsewhere. Alert status contributes exactly
zero and is retained in the table for completeness, because an attribute that survives the
admissibility rules and then turns out to be uninformative is a different finding from one that
was excluded.}
\label{fig:shap}
\end{figure}

The admissible group breach-rate encoding is the strongest single contributor at \num{0.6559},
followed by the assignment delay at \num{0.4909} and the group volume encoding at \num{0.1778}. The
ordering is consistent with the feature ladder: the encodings add little \emph{marginal} AUC at
rung 5 because the configuration-item attributes at rung 4 already carry much of the same
information, but conditional on being present they carry the largest per-case attribution.
Marginal contribution and attribution answer different questions and it is worth saying so
explicitly, because a reader who expects them to agree will read one of the two tables as an
error.

Two attributions deserve comment. \texttt{Alert\_Status} contributes exactly \num{0.0}, which
is consistent with the field being uniformly \texttt{closed} across almost the whole log; it
passes the admissibility rules and is simply uninformative. And \texttt{Assignment\_Delay} ranks
third despite the corruption in its upper tail documented in the Limitations, which suggests the
signal it carries is concentrated in the well-formed short-delay range.

Attribution is reported and is not validated against human decisions. SHAP has known
impossibility results \citep{bilodeau2024impossibility} and sensitivity to feature
representation \citep{hwang2025shapbased}, and whether explanations improve human decisions is
unsettled \citep{haag2026explanations,chae2025explanation}. No claim is made here that these
attributions improve operator performance; they are reported as a description of the model, not
as an intervention.


% =====================================================================================
%  Results for the two scoring formulations. Every number here is written by
%  code/formulations/fit_bpi2014.py into code/results_formulations/.
% =====================================================================================

\section{Results: the scoring formulations}
\label{sec:formulation-results}

Section~\ref{sec:results-predictive} measured a constraint. This section measures the two
formulations built under it. The rule for both is the same and is stated once: every constant
is estimated on the earliest \SI{80}{\percent} of the timeline by incident open time, and every
number reported as an evaluation is computed on the remaining \SI{20}{\percent}, which the
fitting never saw. The working set is the \num{39449} incidents and \num{87} groups of
Section~\ref{sec:two-sets}, of which \num{31559} form the training period.

\subsection{What is fitted, and on how much}

Table~\ref{tab:fitted} lists every estimated constant in F1 and F2 with its interval and its
sample size. Nothing in either formulation is a number chosen because it looked reasonable.

\begin{table}[htbp]
\centering
\caption{Every fitted constant in the two formulations. Intervals are \SI{95}{\percent}. All
estimation uses training-period cases only.}
\label{tab:fitted}
\small
%% GENERATED by code/figures/make_tables.py. Do not edit by hand.
%% The caption, label and float live in the manuscript; only the numbers are here.
\begin{tabular}{lL{5.1cm}rl}
\toprule
Symbol & Quantity & Estimate & Interval / support \\
\midrule
$\hat{\lambda}$ & Freshness decay rate, per hour & 0.0174 & [0.0161, 0.0188], $n=\num{31394}$ \\
 & \quad implied half-life (hours) & 39.9 & derived from $\hat{\lambda}$ \\
$\hat{k}$ & Urgency ramp steepness & 0.0228 & $n=\num{31394}$ \\
$\hat{\tau}_0$ & Urgency ramp midpoint, hours & $-12.03$ & $n=\num{31394}$ \\
$\varepsilon$ & Floor applied before any $p\le0$ aggregation & $10^{-3}$ & fixed, reported as a parameter \\
$\alpha$ & Target-encoding smoothing prior & 20 & fixed \\
$\hat{\theta}$ & F2 escalation cut point & 0.5250 & tuned on train, $F_1=\num{0.5134}$ \\
\midrule
 & Cosine-expertise cells evaluated & \num{686430} & $\num{7890}\times\num{87}$ \\
 & \quad of which exactly zero & \num{206214} & \SI{30.04}{\percent} \\
 & Delay rows above the \num{168}\,h cap & \num{171} & \SI{0.4}{\percent} of \num{39449}, a modelling choice \\
\bottomrule
\end{tabular}


\end{table}

Two of these deserve comment rather than a row.

The freshness decay was the first thing to break, and how it broke is worth recording,
because the first explanation we wrote for it was wrong. Fitted across the assignment delay
column as it was then loaded, the maximum-likelihood estimate of $\lambda$ came back exactly
zero with an undefined confidence interval. The delay column showed a median of \num{16.2} hours
against a 90th percentile of \num{5684} hours, and we recorded that as a corrupted upper tail in
the log's first-assignment timestamp.

It was not the log. The loader was subtracting a month-first incident timestamp from a
day-first activity timestamp, so every delay carried a date-convention error of up to eleven
months, and the tail was ours. Parsed correctly the same column has a median of \num{0.49} hours, a
90th percentile of \num{16.57} and a 99th of \num{92.52}. The \num{168}-hour cap now excludes
\num{171} rows of \num{39449}, \SI{0.4}{\percent}, rather than a third of the variable.

The cap therefore survives as a modelling choice and not as a data repair, and it is reported
as one. A wait of more than a week to first assignment is a different operational regime from
the speed-to-lead window $s_5$ is meant to score, so those cases are held out of the fit. The
consequence is measured rather than waved past: the uncapped control fit is now perfectly
well formed at $\hat{\lambda} = \num{0.0097}\,{\rm h}^{-1}$,
interval $[\num{0.0088}, \num{0.0107}]$, so capping
does not rescue a degenerate fit, it nearly doubles the estimated decay rate to
$\num{0.0174}$. Both are reported. Reading $s_5$ off the capped fit assumes the
organisation cares about the first week, which is a statement about the process and not about
the data.

The urgency midpoint $\hat{\tau}_0 = \num{-12.03}$ hours is negative, which means the fitted
logistic ramp is centred before the observed range begins and $s_6$ is therefore monotonically
decreasing across the whole range rather than sigmoidal within it. The shape is real and it is
reported as fitted. It says that on this log urgency has no inflection point inside the
observable window: cases do not become suddenly urgent at some threshold, they become steadily
less recoverable from the moment they arrive.

\subsection{The criterion decision matrix}
\label{sec:decision-matrix}

Table~\ref{tab:criteria} gives the spread of each criterion over the evaluation period. It is
the input to CRITIC, and it explains most of what the weighting analysis then finds.

\begin{table}[htbp]
\centering
\caption{Criterion statistics over the evaluation-period decision matrix. Conflict is the CRITIC
term $\sum_\ell (1-\rho_{k\ell})$. Exact zeros matter because they are what force the floor in
Equation~\ref{eq:powermean}.}
\label{tab:criteria}
\small
%% GENERATED by code/figures/make_tables.py. Do not edit by hand.
%% The caption, label and float live in the manuscript; only the numbers are here.
\begin{tabular}{llrrrrr}
\toprule
 & Criterion & Mean & SD & Min & Exact zeros & Conflict \\
\midrule
$s_1$  & Outcome propensity       & 0.5036 & 0.2958 & 0.0003 & 0 & 5.397 \\
$s_2$  & Service level survival   & 0.5629 & 0.2904 & 0.0000 & \num{16352} & 5.349 \\
$s_3$  & Expertise match          & 0.3608 & 0.4225 & 0.0000 & \num{109012} & 5.996 \\
$s_4$  & Workload headroom        & 0.3301 & $\num{9.8e-06}$ & 0.3301 & 0 & 6.063 \\
$s_5$  & Freshness                & 0.9470 & 0.1289 & 0.0600 & 0 & 6.529 \\
$s_6$  & Urgency                  & 0.4169 & 0.0659 & 0.3676 & 0 & 6.960 \\
$s_7$  & Importance               & 0.1984 & 0.1690 & 0.0010 & 0 & 5.576 \\
\bottomrule
\end{tabular}


\end{table}

\begin{figure}[htbp]
\centering
\includegraphics[width=0.94\textwidth]{fig06_criterion_spread_and_zeros.pdf}
\caption{Criterion distributions and the exact-zero counts that drive the degeneracy in
Equation~\ref{eq:powermean}.}
\label{fig:criteria}
\end{figure}

Three readings, one of them a warning about our own method.

Expertise is bimodal, not continuous. Its standard deviation of \num{0.4225} is the largest in
the table and larger than its mean, which is the signature of a variable that is mostly at zero
or mostly near one. \num{109012} of \num{348000} cells in this decision matrix are exactly
zero because the group has no training history in that incident's category, and the same
measurement over the full evaluation matrix in Table~\ref{tab:fitted} agrees to within a point.
This is the fact that forces $\varepsilon = 10^{-3}$: without it, every $p \le 0$ aggregation would
return zero for roughly a third of all candidate pairs, and the score would silently become a
filter rather than a ranking.

Workload headroom is almost constant at the batch level, standard deviation
$\num{9.8e-06}$. Within a single arrival batch of 40 incidents against \num{87} groups, no assignment
moves Jain's index far. The criterion does real work when accumulated across the whole
optimisation horizon, which is where the objective $f_3$ measures it, but as a per-pair score it
carries almost no discriminating information, and any weighting scheme that reads spread should
be expected to down-weight it.

Freshness sits near its ceiling, mean \num{0.9470}, because most cases in this log are assigned
quickly. A criterion that is near-saturated for most of the data can still be decisive for the
tail, and the exponential form makes it so, but its \emph{average} contribution is small.

\subsection{Weighting, and whether it matters}
\label{sec:weighting-results}

Table~\ref{tab:weights} gives the three weight vectors. They disagree sharply, which is the
expected result: objective weighting methods are known not to be interchangeable
\citep{mukhametzyanov2021specific,ariyanti2025comparative}, to vary in stability under
perturbation \citep{linh2025evaluating,paradowski2025novel} and to admit rank reversal
\citep{li2026trend}.

\begin{table}[htbp]
\centering
\caption{Weights under the three schemes. Entropy weighting assigns effectively zero to
workload headroom, which is the correct response to a criterion with a standard deviation of
$\num{9.8e-06}$ and an illustration of how differently the two schemes read the same matrix.}
\label{tab:weights}
%% GENERATED by code/figures/make_tables.py. Do not edit by hand.
%% The caption, label and float live in the manuscript; only the numbers are here.
\begin{tabular}{lrrr}
\toprule
Criterion & CRITIC & Entropy & Equal \\
\midrule
s1 outcome     & 0.1628 & 0.1283 & 0.1429 \\
s2 sla survival & 0.1589 & 0.1032 & 0.1429 \\
s3 expertise   & 0.2582 & 0.4687 & 0.1429 \\
s4 workload    & 0.1223 & 0.0000 & 0.1429 \\
s5 freshness   & 0.0913 & 0.0071 & 0.1429 \\
s6 urgency     & 0.0783 & 0.0065 & 0.1429 \\
s7 importance  & 0.1283 & 0.2861 & 0.1429 \\
\bottomrule
\end{tabular}


\end{table}

\begin{figure}[htbp]
\centering
\includegraphics[width=0.94\textwidth]{fig05_weights_and_rank_stability.pdf}
\caption{The three weight vectors and the rank correlations between the orderings they induce.}
\label{fig:weights}
\end{figure}

The disagreement between the vectors is large. Entropy puts \num{0.4687} on expertise and
effectively nothing on workload; CRITIC spreads from \num{0.0783} to \num{0.2582}. A reader is
entitled to ask which one the results depend on, and the honest way to answer is to measure it
rather than to argue for a favourite.

The induced assignment orderings agree closely. Spearman $\rho$ is \num{0.959} between CRITIC
and uniform weighting, \num{0.946} between CRITIC and entropy, and \num{0.856} between entropy
and uniform, all at $p < 10^{-10}$. So the weighting choice is not load-bearing for these
results. That is a stronger statement than a defence of any single scheme, and it is
falsifiable in a way that a defence is not: had the correlations come back at \num{0.4}, the
paper would have had to report that the conclusions depend on a modelling choice with no
principled resolution.

We also state why CRITIC alone would have been unsafe even had it agreed. Two of the seven
criteria, $s_1$ and $s_2$, are model outputs. Their spread reflects how the classifier was
calibrated, and isotonic calibration compresses variance, so a CRITIC weight computed on them
measures a property of our pipeline rather than a property of the decision problem. Publishing
that number alone, labelled data-driven, would have laundered an artifact into a result.

\subsection{The compensation exponent}
\label{sec:psweep}

The exponent $p$ in Equation~\ref{eq:powermean} is what turns a qualitative preference into a
measured one. Table~\ref{tab:psweep} sweeps it.

\begin{table}[htbp]
\centering
\caption{Assignment score against the compensation exponent, evaluation period. The floored
column records whether $\varepsilon$ was engaged, which happens for every $p \le 0$ because
\SI{31.3}{\percent} of expertise cells in this matrix are exactly zero.}
\label{tab:psweep}
%% GENERATED by code/figures/make_tables.py. Do not edit by hand.
%% The caption, label and float live in the manuscript; only the numbers are here.
\begin{tabular}{rrrrr}
\toprule
$p$ & Mean score & Min & Max & Exact zeros \\
\midrule
-4.0 & 0.0949 & 0.0012 & 0.5190 & 0 \\
-2.0 & 0.1135 & 0.0014 & 0.6361 & 0 \\
-1.0 & 0.1336 & 0.0018 & 0.7114 & 0 \\
-0.5 & 0.1565 & 0.0031 & 0.7488 & 0 \\
0.0 & 0.2390 & 0.0116 & 0.7828 & 0 \\
0.2 & 0.2918 & 0.0209 & 0.7982 & 0 \\
0.5 & 0.3631 & 0.0558 & 0.8123 & 0 \\
1.0 & 0.4495 & 0.1387 & 0.8371 & 0 \\
2.0 & 0.5494 & 0.2553 & 0.8743 & 0 \\
4.0 & 0.6659 & 0.3602 & 0.9162 & 0 \\
8.0 & 0.7787 & 0.4480 & 0.9504 & 0 \\
\bottomrule
\end{tabular}


\end{table}

\begin{figure}[htbp]
\centering
\includegraphics[width=0.86\textwidth]{fig04_power_mean_sensitivity.pdf}
\caption{Assignment score against the compensation exponent. The curve is monotone in $p$, as
the power-mean inequality requires, which is a useful implementation check as well as a
property.}
\label{fig:psweep}
\end{figure}

The mean score rises monotonically from \num{0.0949} at $p=-4$ to \num{0.7787} at $p=8$, which
is exactly what the power-mean inequality requires and therefore doubles as a correctness test
on the implementation. The operational content is in the minimum column. At $p=1$ the worst
candidate pair in the evaluation set still scores \num{0.1387}, because a strong criterion can
compensate a weak one; at $p=-4$ it scores \num{0.0012}, because it cannot. An organisation
that never wants an incident routed to a group with no relevant history should choose a
negative $p$ and accept that the floor $\varepsilon$ then determines what "no history" is worth.
An organisation willing to trade expertise for availability should choose $p \ge 1$. The
formulation does not decide this; it makes the decision explicit and reports what it costs.

\subsection{F2: deriving breach beats predicting it}
\label{sec:hazard-results}

Table~\ref{tab:hazard} reports the service level formulation. The comparison is between two
models of the same event on the same data, one of which is structurally able to observe its own
label and one of which is not.

\begin{table}[htbp]
\centering
\caption{Service level risk. The person-period expansion turns \num{39449} cases into
\num{161720} rows over ten elapsed-time bins. Every AUC is on the held-out evaluation period.}
\label{tab:hazard}
%% GENERATED by code/figures/make_tables.py. Do not edit by hand.
%% The caption, label and float live in the manuscript; only the numbers are here.
\begin{tabular}{lrl}
\toprule
Model & Held-out AUC & What it estimates \\
\midrule
Direct breach classifier & 0.7879 & $P(\text{breach} \mid \mathbf{x})$ at intake \\
Resolution hazard, bin level & 0.7640 & $h(t \mid \mathbf{x})$ per person-period row \\
Breach derived from resolution, at intake & \textbf{0.8373} & $S(\tau\mid\mathbf{x})/S(t\mid\mathbf{x})$ \\
\midrule
Breach hazard with elapsed time (rejected) & 0.9542 & reads its own label \\
\bottomrule
\end{tabular}


\end{table}

The derived quantity is better than the direct classifier by \num{0.0494} AUC. This is the
result we did not expect and it inverts the usual framing. Admissibility is normally described
as a constraint that costs performance, and Section~\ref{sec:cost} measures exactly that cost.
Here, choosing the object that can actually be estimated \emph{improved} estimation, because the
resolution hazard is a genuine random quantity given elapsed time whereas breach, past the
threshold, is not.

The rejected row is included deliberately. A breach hazard given elapsed time and priority
scores \num{0.9542}, which is by a wide margin the best number in this paper and is worthless.
The reason is a property of the label, and we test it without reconstructing the label's own
threshold, since rederiving a formula only proves it gives the same answer twice. Within every
priority class the longest non-breaching case is shorter than the shortest breaching one, so
breach is a step function of handle time: of the \num{16111} cases at or past that cut, the
measured breach rate is exactly \num{1.0000}. The model is reading the definition of its
label rather than predicting an event. It is reported because a paper about claims that fail
verification should show its own worked example of the failure, and because a reviewer
encountering a number of this size elsewhere in this literature now has a specific mechanism
to test for.

This row is also our own smallest cautionary tale. It stood in the manuscript as a hardcoded
constant for several drafts, carried forward from an earlier run of a dataset we had since
replaced, in the one section of the paper whose subject is numbers that nobody re-derived. It
is now measured in the same run as every other row of Table~\ref{tab:hazard}, on the same
split, with the same features, changing only the event definition.

Two checks are run on the derived quantity rather than assumed. Survival $S(t\mid\mathbf{x})$ is
verified monotone non-increasing over the bin sequence, which it must be by construction and
which fails loudly if the hazard is mis-indexed. And the per-bin resolution rate is
\num{0.2430}, comfortably away from both zero and one, so the hazard model is estimating a
balanced event rather than a rare one.

\subsection{Workload, reported two ways}

Over the \num{87} groups of the evaluation period the load vector has Jain index \num{0.3301},
standard deviation \num{516.7} and coefficient of variation \num{1.424}. Both the bounded and
the unbounded measure are reported for every policy in Section~\ref{sec:assignment-results} and
neither is substituted for the other, because the proposed method loses on workload spread and
replacing the metric it loses on after seeing the result would be fitting the metric to the
outcome.


% =====================================================================================
%  Multi-objective assignment results. All numbers from code/phase4/results/*_fullv3.*
%  1,225 runs: 13 policies, 5 batches, 30 seeded repeats, 150 generations, population 60.
% =====================================================================================

\subsection{Multi-objective assignment}
\label{sec:assignment-results}

The protocol runs five operational baselines, standard NSGA-II, NSGA-III as a many-objective
control, the proposed PAA-NSGA-II, its three single-component ablations, and two post-hoc
exploratory combinations. Five arrival batches, 30 seeded repeats each, 150 generations,
population 60. All \num{1225} runs completed with \textbf{zero capacity violations}, so the repair
operator satisfies constraint C2 throughout.

\begin{table}[htbp]
\centering
\caption{Assignment policies on BPI Challenge 2014. Hypervolume is computed against a
reference point fixed as the nadir across all variants and all runs within a batch. Objectives
$f_1$ and $f_4$ are maximised, $f_2$ and $f_3$ minimised. Jain's index is reported alongside
the $f_3$ standard deviation and is not substituted for it.}
\label{tab:assignment}
\small
%% GENERATED by code/figures/make_tables.py. Do not edit by hand.
%% The caption, label and float live in the manuscript; only the numbers are here.
\begin{tabular}{lrrrrr}
\toprule
Policy & Hypervolume & $f_1$ FTR & $f_2$ breach & $f_3$ workload SD & $f_4$ expertise \\
\midrule
Random                         & -- & 0.4890 & 0.4351 & 0.7807 & 0.3562 \\
Round robin                    & -- & 0.4592 & 0.4510 & 1.1796 & 0.3449 \\
FIFO                           & -- & 0.4500 & 0.5523 & 2.5554 & 0.4039 \\
Least loaded                   & -- & 0.4473 & 0.4528 & 0.4984 & 0.3876 \\
Greedy best outcome            & -- & 0.6998 & 0.2885 & 1.6501 & 0.7162 \\
Standard NSGA-II               & 0.0695 & 0.6169 & 0.3146 & 1.0352 & 0.6988 \\
NSGA-III                       & 0.0824 & 0.6455 & 0.2792 & 0.9978 & 0.7556 \\
PAA-NSGA-II                    & 0.0762 & 0.6306 & 0.3014 & 1.2552 & 0.7478 \\
PAA without priority init      & 0.0658 & 0.6156 & 0.3076 & 1.0467 & 0.6644 \\
PAA without adaptive crossover & 0.0767 & 0.6325 & 0.3028 & 1.2518 & 0.7594 \\
PAA without SLA mutation       & 0.0767 & 0.6303 & 0.2964 & 1.2414 & 0.7373 \\
PAA with adaptive mutation     & 0.0766 & 0.6337 & 0.3046 & 1.2710 & 0.7657 \\
NSGA-III + priority init       & 0.0850 & 0.6595 & 0.2834 & 1.1679 & 0.7931 \\
\bottomrule
\end{tabular}



\vspace{0.4em}
{\footnotesize \textsuperscript{*}Post-hoc exploratory combination, designed after observing
the ablation. Reported as exploratory and not claimed as a contribution.}
\end{table}

\paragraph{The proposed operators improve on their own base algorithm.}
Against standard NSGA-II over 150 paired runs, PAA-NSGA-II is better on four of five metrics
by paired Wilcoxon signed-rank with Holm-Bonferroni correction: hypervolume
($+\num{0.00666}$, $p = \num{1.5e-19}$), first-time-right ($+\num{0.01367}$,
$p = \num{1.3e-5}$), breach ($-\num{0.01322}$, $p = \num{1.5e-4}$) and expertise match
($+\num{0.04898}$, $p = \num{9.9e-7}$).

\paragraph{The base algorithm matters more than the modifications.}
NSGA-III, included only as a many-objective control, reaches hypervolume \num{0.0824} against
\num{0.0762} for the proposed method, a paired difference of $\num{-0.00625}$ at
$p = \num{1.5e-23}$. It does so with a front of about \num{18} solutions against \num{60}, so it
attains higher hypervolume from under a third as many points. With five objectives this is the
regime for which reference-point selection was introduced
\citep{debjain2014evolutionary,zheng2024runtime}, and the result is consistent with that
literature rather than in tension with it.

\paragraph{One of three components does the work.}
Table~\ref{tab:ablation} isolates each modification. Removing priority-aware initialisation
costs \num{0.0104} hypervolume, so that component is doing real work. The other two are
not: removing adaptive crossover \emph{raises} hypervolume by \num{0.0005} and removing
SLA-aware mutation raises it by \num{0.0006}, both small enough to sit inside run-to-run
noise but neither of them a contribution. Reporting the three as a bundle would have credited
them equally, and that would have been wrong in both directions.

\begin{table}[htbp]
\centering
\caption{Ablation. Change in hypervolume when a single component is removed from the proposed
method. A negative delta means the component contributes.}
\label{tab:ablation}
%% GENERATED by code/figures/make_tables.py. Do not edit by hand.
%% The caption, label and float live in the manuscript; only the numbers are here.
\begin{tabular}{lrr}
\toprule
Component removed & Hypervolume without & $\Delta$ \\
\midrule
Priority-aware initialisation  & 0.0658 & $-0.0104$ \\
Adaptive crossover             & 0.0767 & $+0.0005$ \\
SLA-aware mutation             & 0.0767 & $+0.0006$ \\
\midrule
None (full method)             & 0.0762 & \\
\bottomrule
\end{tabular}


\end{table}

\begin{figure}[htbp]
\centering
\includegraphics[width=0.72\textwidth]{fig08_ablation.pdf}
\caption{Component ablation. Removing priority-aware initialisation costs \num{0.0104}
hypervolume; removing either of the other two changes it by less than \num{0.001} and in the
opposite direction, which is inside run-to-run noise. Reporting the three as a bundle
would have credited them equally.}
\label{fig:ablation}
\end{figure}

\paragraph{The best configuration is not a significant improvement.}
Placing the one modification that contributes onto the base algorithm that wins gives the
highest hypervolume observed, \num{0.0850}. Against plain NSGA-III the paired difference is
$\num{+0.00254}$ at $p = \num{6.8e-11}$, which does clear $\alpha = \num{0.05}$. It is still
reported as exploratory and is still not claimed, because it was designed after reading the
ablation table on this same data. A post-hoc combination that clears a threshold on the
data that suggested it has not been tested; it has been fitted. Claiming it would be the
same error this paper is about, committed one level up.

\paragraph{A real trade-off, visible in two independent metrics.}
The proposed method loses significantly on workload spread ($+\num{0.22006}$,
$p = \num{1.6e-21}$). Jain's fairness index agrees: \num{0.1207} against
\num{0.1716} for standard NSGA-II. The two metrics were computed independently and could have disagreed. The
method concentrates work on capable groups to gain outcome quality, which is a trade-off an
operator may or may not want. The exponent $p$ in Equation~\ref{eq:powermean} is the parameter
through which that preference is expressed rather than assumed.

\begin{figure}[htbp]
\centering
\includegraphics[width=0.78\textwidth]{fig07_hypervolume_by_variant.pdf}
\caption{Search quality across algorithm variants. The proposed method improves on its own
base algorithm and is in turn outperformed by an off-the-shelf many-objective control.}
\label{fig:hypervolume}
\end{figure}


% =====================================================================================
%  Complexity and measured runtime. Numbers from code/results_evaluation/runtime.csv
%  and complexity.csv. Timing from an instrumented run: 160 solves, 20 per variant.
% =====================================================================================

\subsection{Complexity and measured cost}
\label{sec:complexity}

\paragraph{Asymptotic cost.}
Write $n$ for the number of work items in an arrival batch, $m$ for the number of eligible
groups, $N$ for population size, $M$ for the number of objectives and $H$ for the number of
reference directions. The operational baselines are single-pass: random, round robin and FIFO
are $O(n)$, least-loaded is $O(n \log m)$ with a heap over group loads, and greedy is $O(nm)$
because it scans every eligible group for every item. The evolutionary methods are dominated
per generation by non-dominated sorting at $O(M N^2)$ \citep{deb2002fast}; NSGA-III adds
association to reference directions, giving $O(M N^2 + N H)$ \citep{debjain2014evolutionary}.
The repair operator is $O(nm)$ in the worst case and is bounded at six passes, so it does not
change the order.

The proposed modifications do not alter the asymptotic cost. Priority-aware initialisation is
a one-off $O(nm)$ greedy construction over a fraction of the initial population. Adaptive
crossover and the mutation schedule each add an $O(N n)$ diversity measurement per generation,
dominated by the $O(M N^2)$ sorting term.

\paragraph{Measured cost.}
Asymptotic equality does not imply equal wall time, so Table~\ref{tab:runtime} reports
measured per-solve time from an instrumented run of 160 solves, 20 per variant, at 150
generations and population 60 on CPU with no GPU. Timing is captured at the point the search
executes rather than derived from log timestamps.

\begin{table}[htbp]
\centering
\caption{Measured wall time per solve. 20 solves per variant, 150 generations, population 60,
single CPU. Percentages are relative to standard NSGA-II.}
\label{tab:runtime}
%% GENERATED by code/figures/make_tables.py. Do not edit by hand.
%% The caption, label and float live in the manuscript; only the numbers are here.
\begin{tabular}{lrrrr}
\toprule
Policy & Mean s/solve & SD & Solves & Percent vs NSGA-II \\
\midrule
Standard NSGA-II               & 0.89 & 0.02 & 20 & 0.0 \\
NSGA-III                       & 0.99 & 0.01 & 20 & 12.2 \\
PAA-NSGA-II                    & 0.96 & 0.01 & 20 & 8.1 \\
PAA without priority init      & 0.96 & 0.01 & 20 & 7.8 \\
PAA without adaptive crossover & 0.91 & 0.02 & 20 & 2.4 \\
PAA without SLA mutation       & 0.98 & 0.02 & 20 & 10.1 \\
PAA with adaptive mutation     & 0.92 & 0.01 & 20 & 3.5 \\
NSGA-III + priority init       & 1.01 & 0.02 & 20 & 14.3 \\
\bottomrule
\end{tabular}


\end{table}

\begin{figure}[htbp]
\centering
\includegraphics[width=0.72\textwidth]{fig09_runtime.pdf}
\caption{Measured wall time per solve, 20 solves per variant, single CPU. Every evolutionary
variant lies within \SI{14}{\percent} of standard NSGA-II, so on this problem the choice
between them is not a cost decision.}
\label{fig:runtime}
\end{figure}

Three observations follow.

The spread is narrow. Every evolutionary variant lies within \SI{14}{\percent} of standard
NSGA-II, and at \num{0.89} to \num{1.01} seconds per batch of 40 items, with a standard
deviation under \num{0.02} across twenty solves each, all are far inside the latency budget of an
assignment decision a human would otherwise take in seconds. These timings were measured on
an otherwise idle machine. An earlier set, taken while four other jobs were running, read two
to three times higher and with twenty times the variance, and is recorded here because a
runtime measured under contention is not a runtime. These timings were measured on a workstation that was not
otherwise idle, so they are reported as indicative of the ordering and not as benchmarks.

Priority-aware initialisation buys search quality and is paid for in time. Adding it to
NSGA-III raises hypervolume from \num{0.0837} to \num{0.0862}, and raises mean solve time
from \num{1.80} to \num{3.19} seconds, which is \SI{77}{\percent} longer.

An earlier draft reported this as a component that pays for itself twice, raising quality
and cutting time together. That was measured before the hypervolume reference point was
fixed and before the timing instrumentation was corrected, and it does not survive either.
Constructing a priority-aware starting population costs real work and the run does not win
it back. The trade is still favourable at this scale, because the extra \num{1.38} seconds sits
well inside the budget of a decision a human would take in seconds, but it is a trade and
it is reported as one.

The two modifications that do not help do not cost much either. Removing adaptive crossover or
SLA-aware mutation saves roughly \SI{5}{\percent} of solve time while changing hypervolume by
less than \num{0.001}. They are close to free and close to useless, which is the honest
summary and the reason the ablation is reported per component.


% =====================================================================================
%  Discussion, threats to validity, limitations. Included by paper.tex after Results.
% =====================================================================================

\section{Discussion}

\subsection{A fifth of the signal was not there}

The headline cost figure is \num{0.0792} AUC, or \SI{19.6}{\percent} of everything the
unconstrained model achieved above chance. That is not a rounding correction. A practitioner
deploying the unconstrained model would find a fifth of its apparent discrimination absent
at run time, because the attributes carrying it are recorded only after the decision has been
taken. The same measurement under an \num{80}/\num{20} split differs by less than
\num{0.002} AUC, so the estimate is a property of the data rather than of the split.

This is consistent in direction with reported effects elsewhere, including the
\num{15.07} point accuracy drop \citet{mishra2026taxonomy} measure on build prediction and
the example-leakage rates above \SI{80}{\percent} that \citet{abb2024discussion} find on
standard process logs. The contribution here is that the cost and the benefit are measured on
the same log under one protocol, so neither can be quoted without the other.

\subsection{Admissibility is not the same as parsimony}

A natural response to a leakage finding is to use fewer attributes. The ladder shows that
would be the wrong lesson. Discrimination rises monotonically from \num{0.5100} at three
admissible attributes to \num{0.8270} at seventeen, a gain of \num{0.2560} AUC, reproduced to
within \num{0.0053} across split ratios. Restricting to a small attribute set costs far more
than enforcing admissibility does.

The distinction matters because the two are easily conflated. Attributes should be dropped
because they are unavailable at decision time, not because they are numerous. Nine attributes
are excluded here, each with its reason recorded; the remaining seventeen are kept because
each is observable when the assignment is made.

\subsection{Choosing the admissible target can improve accuracy, not only defensibility}

The service level result is the one we did not expect. Breach cannot be modelled directly as
a survival event on this log without the model reading its own label, because breach is
defined by duration exceeding a per-priority threshold and the measured breach rate past that
threshold is exactly \num{1.0000}. An early version scored \num{0.9418} AUC on precisely that
artifact.

Modelling resolution instead, and deriving breach as the conditional survival ratio
$S(\tau \mid x) / S(t \mid x)$, produces \num{0.8373} AUC at intake against \num{0.7879} for a
direct classifier. The admissible formulation is better by \num{0.0494} AUC. Admissibility is
usually framed as a constraint that costs performance. Here, choosing the object that can
actually be estimated improved it, because the resolution hazard is a genuine random quantity
given elapsed time while breach, past the threshold, is not.

\subsection{The weighting objection answers itself}

Objective weighting methods are known to disagree, to be unstable under perturbation, and to
admit rank reversal \citep{mukhametzyanov2021specific,linh2025evaluating,li2026trend}. CRITIC
is additionally unsuitable in isolation here, because two of the seven criteria are model
outputs whose spread reflects calibration rather than importance, and isotonic calibration
compresses variance. Reporting CRITIC alone would convert a modelling artifact into a weight
and label it data-driven.

Reporting three schemes resolves this empirically rather than rhetorically. The induced
assignment orderings agree at Spearman $\rho = \num{0.959}$ between CRITIC and uniform
weighting and $\rho = \num{0.946}$ between CRITIC and entropy. The weighting choice is
therefore not load-bearing for these results, which is a stronger statement than defending any
one scheme.

\subsection{The modifications work, but the base algorithm matters more}

The proposed operator set does improve on standard NSGA-II, significantly and on four of five
metrics (Section~\ref{sec:assignment-results}). It is then outperformed by NSGA-III, an
off-the-shelf control included only as a many-objective reference, which reaches higher
hypervolume from a front roughly a third the size. Two considerations make this informative
rather than disappointing.

First, five objectives place this problem in the many-objective regime where dominance-based
selection is known to degrade, a limitation characterised when NSGA-III was introduced
\citep{debjain2014evolutionary} and still under active theoretical study
\citep{zheng2024runtime,doerr2026difficulties}. A result in which a reference-point method
outperforms a dominance-based one on five objectives agrees with that theory.

Second, \citet{wang2025rethinking} argue that much claimed novelty in metaheuristic research
does not survive controlled comparison. An ablation that isolates each proposed modification
is the appropriate response to that argument, and it is why the ablation is reported per
component rather than only in aggregate.

\subsection{What the cost looks like in cases rather than in AUC}

A \num{0.0792} AUC loss is an abstraction. The confusion matrices in
Section~\ref{sec:confusion} put it in units an operations manager can act on. On the same
\num{13982} held-out incidents, the unconstrained model catches \num{4295} of the \num{5136}
breaches and the admissible one catches \num{3795}. The \num{500} cases in that difference are
not a modelling subtlety. They are the incidents an escalation queue built on the
unconstrained model would appear to catch and would not.

The direction of the error also changes. The unconstrained model is roughly symmetric,
\SI{82.3}{\percent} specificity against \SI{83.6}{\percent} recall. The admissible model is
not: \SI{74.6}{\percent} specificity against \SI{73.9}{\percent} recall. Removing the
post-hoc attributes removes signal about the long-running cases specifically, because those
attributes counted the events a long-running case accumulates. What is left behind
discriminates the ordinary case well and the pathological case poorly, which is the harder
half of the problem and the half a service desk actually cares about.

\subsection{The model is data-limited, not capacity-limited}

The learning curve in Section~\ref{sec:learning} separates two explanations for
\num{0.8270} that are often conflated. Training AUC is \num{0.9999} at \num{1631} incidents,
so the learner has ample capacity; the ceiling is not the model class. The generalisation gap
falls monotonically from \num{0.2234} to \num{0.0563} across a twentyfold increase in data
and is still falling, and held-out AUC is still rising at the largest training size available.

Two consequences follow, and the second is a warning against the obvious fix. More incidents
from this desk would help, so the practical route to a better admissible model is more history
rather than a better algorithm, which is consistent with the four learners landing within
\num{0.0158} AUC of each other. And heavier regularisation is the wrong first response: it
would close the gap by moving the training curve down, not by moving the held-out curve up.
The gap is not evidence that the model is over-parameterised. It is evidence that
\num{32623} incidents is not yet enough for a 17-attribute problem with this much label noise.

\subsection{Ranking and operating point fail separately}

The chronological result is usually reported as one number, and Section~\ref{sec:temporal}
shows it is two failures with different remedies. Ranking degrades from \num{0.8258} to
\num{0.8074} AUC and no threshold recovers that. The operating point degrades as well, from
$F_1 = \num{0.6790}$ on the random split to $F_1 = \num{0.6087}$ at the same default cut point,
and tuning the cut point on training-period data alone recovers it to \num{0.6157} without
retraining anything. Recall is where the operating point actually moves: \SI{67.7}{\percent} at the
default against \SI{81.0}{\percent} tuned.

Conflating the two produces the wrong operational conclusion in both directions. A team that
sees only the fall in $F_1$ will retrain a model that did not need retraining. A team that
sees only the tuned $F_1$ will not learn that its threshold needs re-tuning as the process
drifts, which on this log it does within a year: median handle time falls from \num{4.28} to
\num{2.79} hours between the training and evaluation periods.

\subsection{Six defects in our own work, and what each one cost}
\label{sec:harness-defects}

Six errors were found in this work and corrected before any number reported here was
produced. All six are recorded because a paper about claims that do not survive verification
should hold its own instrument to the same standard, and because an earlier version of this
section did not: it named three and gave the other three either a sentence elsewhere or
nothing.

That omission was itself the more serious error, so the ordering here is deliberate. The
first three below are defects in what was measured. They change what the paper is entitled
to claim. The last three are defects in the instrument. They change how reliably a number
can be reproduced, which matters, but a reproducible measurement of the wrong thing is
still the wrong thing.

\paragraph{Validity 1: the incident file was a partial copy.} The working copy of
\texttt{Detail\_Incident.csv} held \num{31238} rows against the canonical
\num{46606}, a third of the log missing, and it was a prefix rather than a sample. The file is
ordered by incident ID, which tracks open time, so a prefix is close to a truncation in time.
That is why it cost so much more on the chronological arm than on the random one:
\num{0.0307} against \num{0.0043}, measured by reintroducing it alone. The canonical file was
re-downloaded from the 4TU DOI, and the digest of every raw input is now recorded in
\texttt{code/input\_checksums.json} and checked by \texttt{evaluation/verify\_input\_data.py};
the incident header is MD5 \texttt{b0608464312c135da289d99911735a5d}. A digest stored on the first day
would have caught this on the first run, and every artifact computed from the partial file
was computed correctly from the wrong input, which is the failure mode a checksum exists
for.

\paragraph{Validity 2: a feature of ours violated our own Rule 1.} The standing-queue
count at a case's arrival was computed over cases that eventually resolved, excluding the
\num{1780} that never do. Whether a case will ever resolve is not knowable at
$t^{\mathrm{dec}}$, so the feature read the future. It is worth stating plainly that this
is a violation of the rule this paper exists to state, committed by its authors, and that it
was found by applying the rule mechanically rather than by intuition. Reintroducing it alone
raises the admissible AUC by \num{0.0015}, which is the direction a leak should move a
score and is the reason we believe the diagnosis.

\paragraph{Validity 3: Rule 2 was not implemented as Rule 2.} The pipeline refitted the
group encoding inside every cross-validation fold. That prevents an encoding from seeing the
rows in its own evaluation fold, and it is what most careful work in this area does, but it
is not Equation~\ref{eq:rule2}: under a random split a case from 2012 could still be encoded
using cases that resolved in 2014. The paper's contribution is that equation, so for a period
the protocol described was not the protocol that ran. It is now implemented as written, in
\texttt{code/causal\_encoding.py}, checked against an $O(n^2)$ reference implementation to a
maximum absolute difference of zero. Reintroducing the fold refit alone raises the admissible
AUC by \num{0.0043}. The distinction is therefore load-bearing for the claim and not for the
number, and both halves of that sentence are reported.

\paragraph{Instrument 1: the seed was not a seed.} It was derived using Python's built-in string hash, which is salted per process. The
same variant therefore drew a different seed on every invocation, measured at \num{43},
\num{41} and \num{32} across three processes for one variant name. The harness documentation
asserted reproducibility that did not hold. It now uses a stable checksum, and two independent
invocations produce bitwise-identical metrics.

\paragraph{Instrument 2: the hypervolume reference point moved.} It was taken from whichever run completed first, making every
hypervolume value dependent on variant ordering and incomparable between runs with different
variant sets. It is now the nadir across all variants and all runs on a batch, computed before
any hypervolume is taken.

\paragraph{Instrument 3: the thread count was never pinned.} This one most directly
threatens what the paper claims about verification. Every number here is
printed to four decimals and checked against a tolerance of $\num{0.00005}$. That tolerance
only means something if a re-run reproduces the number, and it did not. Running the headline
pipeline with one seed, one dataset and one split, changing nothing but the OpenMP thread
count, moves the admissible AUC by \num{0.0015} (2 gives \num{0.8267}, 4 gives \num{0.8258}, 8 gives \num{0.8273})
and the unconstrained AUC by \num{0.0003}. That is roughly \num{30} times the tolerance,
from a variable that appeared in no configuration file and in no methods section.

The cause is not randomness and no seed can absorb it. Gradient-boosted implementations
reduce per-thread histogram partials in thread order, floating-point addition is not
associative, and the resulting difference in a split point compounds through the trees. The
control confirms it: \num{2} runs at the same thread count are bitwise identical,
so the variation is entirely explained by how the work was divided, not by what the code did
with it. The thread count is now pinned in the harness and recorded in every run log, and the
measurement above is reproduced by \texttt{evaluation/thread\_determinism.py}.

We report it rather than quietly pinning it because the implication generalises past this
paper. A reviewer who re-runs a published gradient-boosted pipeline on different hardware
should not expect the fourth decimal to agree, and a paper that verifies its own numbers to
four decimals without fixing the thread count is verifying them against a moving target. Ours
was. Pinning does not make the numbers portable across CPU architectures either, since vector
width changes the summation order within a thread as well; it removes the one source that
varied on a single machine, which was the source that had been silently moving our own
results between runs.

The first two instrument defects do not affect the predictive results in
Sections~\ref{sec:cost} to \ref{sec:hazard-results}, which come from a separate pipeline. Both
invalidated every earlier assignment run, and those were discarded.

The consequence is worth stating plainly rather than in passing. Under the uncorrected
harness the proposed method appeared to LOSE to standard NSGA-II on hypervolume. Corrected,
it wins, significantly. The defect did not merely add noise: it reversed the sign of the
study's headline comparison. Had the harness not been audited, this paper would have
reported the opposite conclusion with the same apparent statistical confidence.

\subsection{Threats to validity}

\textbf{Construct.} The service level label is derived from handle time against \emph{twice} the
per-priority median rather than from a contractual target, because the log's own
\texttt{Alert Status} field is uniformly \texttt{closed} across \num{46606} of \num{46809} raw
incidents and carries no usable breach signal. Results should be read as conditional on that
operationalisation, and Equation~\ref{eq:label} states it exactly so it can be varied. On the
random-split arms the median is taken over the whole log. That is a property of the label and
not of any feature, no model observes it, and every configuration is scored against the same
fixed label, so it cannot affect the comparison between configurations; the chronological arm
re-derives it from training-period cases alone, where it would otherwise matter.

\textbf{Statistical.} The fold-level Wilcoxon signed-rank test at $n=5$ has a smallest
attainable two-sided $p$-value of \num{0.0625} and therefore cannot reject at
$\alpha=\num{0.05}$ regardless of effect size. Every fold-level comparison in
Section~\ref{sec:cv} returns exactly that floor, which means one learner won on all five folds
and not that the difference is small. Model comparisons are therefore made with 2000-resample
bootstrap intervals on the test set, and the fold-level result is reported for transparency
rather than used as evidence.

\textbf{Internal.} Every parameter is fitted on the earliest \SI{80}{\percent} of the timeline,
and the escalation cut point is tuned on training-period data alone. The chronological holdout
is reported alongside the random split precisely because the gap between them, \num{0.8270}
against \num{0.8085}, is itself informative about drift. That gap was \num{0.1326} before the date-parsing defect in the loader was corrected and is \num{0.0185} after it, so most of what earlier drafts read as distribution shift was the shift in how their own timestamps were parsed.

\textbf{External.} The protocol is measured on five public event logs spanning four domains:
IT service management, consumer lending, acute care, healthcare administration and public
administration. That is enough to say the protocol transfers and enough to say the magnitude
does not. The share of above-chance signal lost to admissibility ranges from \SI{9.7}{\percent}
on Hospital Billing to \SI{69.5}{\percent} on Road Traffic Fines. Any single figure quoted from
one log, including the one this study began with, is a point on that range and not a constant.

Four of the five logs are Dutch or Italian European public releases and all five are historical.
Nothing here speaks to a live deployment, to a non-European process, or to an organisation that
knows its assignment policy is being measured.

\textbf{Conclusion.} Statistical comparisons use paired Wilcoxon signed-rank tests with
Holm-Bonferroni correction across five metrics, and every run is seeded and reproducible.

\section{Limitations}

\textbf{An IT service desk is not a sales pipeline.} This work is motivated by lead
assignment, and the objective analogous to conversion is first-time-right resolution. No
public dataset carries real leads, real individual representative identifiers and real
conversion outcomes together. That absence is a finding rather than an oversight, and it is
why the evaluation is conducted on an operational log whose assignment decisions and outcomes
are real even though its domain is not sales.

\textbf{A parsing defect in an earlier version of this work, and what it cost.} The incident
file records timestamps day-first; an earlier version of this pipeline parsed them month-first,
which silently transposed day and month on the \SI{41.0}{\percent} of rows where both fields
are at most twelve. Two consequences were reported as properties of the data and were not.
The observation window was given as \num{1055} days against a true \num{785}. And the
first-assignment delay appeared corrupted in its upper tail, with a median of \num{16.2}
hours against a 90th percentile of \num{5684}, which led to \SI{35.6}{\percent} of rows being
excluded from the freshness and urgency fits as log artifacts. Under the correct parse that
attribute has a median of \num{0.34} hours and a 90th percentile of \num{15.48} over the
\num{46605} incidents of the predictive matrix, and the exclusion above the \num{168}-hour cap
falls to \SI{0.4}{\percent}. Table~\ref{tab:fitted} quotes \num{0.49} and \num{16.57} for the same
attribute because the formulations are fitted on the \num{39449} incidents of the assignment
working set rather than on all of them; the two populations differ and so do their medians. Nothing was corrupt. The subtraction was mixing two
date conventions. Every number in this paper is computed after that correction, and the
defect is recorded here because a paper about attributes that mislead should say when its
own instrument did.

\paragraph{What the corrections do not account for.} Each of the six was reintroduced on
its own, everything else held corrected, and the pipeline re-run
(\texttt{evaluation/attribution\_run.py}). The date defect explains most of the change on
the random split, \num{0.0389} of \num{0.0399}, and the truncated file explains most of it
on the chronological arm, \num{0.0307} of \num{0.1543}. But all six reintroduced together
give \num{0.7996} and \num{0.7235}, while the figures an earlier version of this work
published were \num{0.7859} and \num{0.6533}. So the reconstruction accounts for about two
thirds of the change on the random split and just over half of it on the chronological arm.
The remaining \num{0.0137} and \num{0.0702} are \emph{unexplained}. We do not attribute
them to the thread count: that effect is measured above at \num{0.0015} and \num{0.0019},
which is forty times too small to cover the chronological residual, and offering it as a
candidate would be exactly the kind of unfalsifiable explanation this paper objects to
elsewhere. The earlier pipeline was rewritten rather than parameterised, so it cannot be
re-run to settle the question, and a reader is entitled to treat the pre-correction figures
as unreproducible rather than as merely superseded.

\paragraph{Whether the corrections are self-serving.} They moved the headline from
\num{0.7859} to \num{0.8258} and the chronological arm from \num{0.6533} to \num{0.8076}, and a
round of bug fixing that improves nearly every number deserves the suspicion it attracts. Two
things answer it, and neither is an assurance. First, the two leaks we removed moved the score
\emph{up} when reintroduced, \num{0.0015} and \num{0.0043}, which is the direction a leak must move it;
had removing a leak raised our score, the diagnosis would have been wrong. Second, and more
to the point, several corrections cost us. The four learners are no longer statistically
distinguishable, where an earlier draft reported two of three gaps excluding zero. Two of
the three proposed algorithm components now \emph{raise} the objective when removed, so
they contribute nothing. A component described as paying for itself twice in fact costs
\SI{77}{\percent} more runtime, and that claim is withdrawn. And the study's most quotable finding,
that a deployed model would recover under a tenth of future breaches at its default
threshold, was an artifact of the date defect: the true figure is \SI{67.7}{\percent}, and the
dramatic version is withdrawn rather than softened. A correction round that only ever helped
would be evidence of selection. This one did not.

\textbf{Counterfactual outcomes are not observed.} A historical log records only the
assignment that happened. Any policy differing from the recorded one is scored by a model
rather than by ground truth, which is the standard objection to optimisation studies on
logged data. We report the fraction of decisions matching the recorded assignment together
with the real observed breach rate on exactly those cases, as the only non-simulated outcome
signal available. Off-policy evaluation would address this properly and is left to future
work.

\textbf{Group-level and individual-level assignment both appear, and the distinction matters.}
BPI 2014 identifies assignment groups rather than individuals, so representative-level features
such as experience are not estimable on it. Three of the four added logs do carry individual
resource identifiers, and the encoding runs at that level there: \num{120} distinct originators
on BPI Challenge 2017 and \num{894} on Hospital Billing. The protocol is unchanged between the
two settings, because Rule 2 is stated over whatever entity the log identifies. What we cannot
say is whether encoding at individual level is better than at team level, since no log here
offers both and a comparison across logs would confound the entity with everything else.

\textbf{Explanation is reported but not validated against human decisions.} SHAP attributions
are computed for the predictive layer, but attribution methods have known impossibility
results \citep{bilodeau2024impossibility} and sensitivity to feature representation
\citep{hwang2025shapbased}, and whether explanations improve human decisions is unsettled
\citep{haag2026explanations,chae2025explanation}. No claim is made that these explanations
improve operator performance.


% The abstract, conclusion and declarations are SHARED with the other driver via
% \input. They used to be duplicated inline in both files and they silently
% diverged: the IEEE build shipped a code-availability URL the other did not, and both
% PDFs went out that way. Never inline them again.
% Shared block. Included by BOTH paper.tex and paper_ieee.tex.
% Do NOT duplicate this text into either driver. The two drivers previously carried
% their own copies of the abstract, conclusion and declarations, and they diverged:
% the IEEE build shipped a code-availability URL that the single-column build did not,
% and both PDFs went out that way on 2026-09-07.
\section{Conclusion}

Decision-time admissibility costs \num{0.0792} AUC on the incident log for the strongest
learner, close to a fifth of all above-chance signal, and the same correction appears for all
four learners with every bootstrap interval excluding zero. Across five logs the cost ranges
from \SI{9.7}{\percent} to \SI{69.5}{\percent} of above-chance signal, so no single figure
from any one log, including this one, should be read as a general constant.

What generalises is the shape rather than the magnitude. Contaminated models across the five
logs occupy a band of \num{0.155} AUC; admissible models on the same logs occupy
more than twice that, \num{0.291}. Attributes recorded after a decision make unlike problems look alike. That is the
practical case for the protocol: not that it improves a model, but that it stops a leaderboard
from hiding which problems are actually solved.

The discipline is nevertheless worth paying for on its own terms. An admissible ladder still
gains \num{0.2542} AUC between eight and seventeen attributes, so the lesson of a leakage
finding is not to model with fewer attributes but to model with observable ones. Choosing the
admissible object rather than the convenient one can improve accuracy outright, as deriving
breach from a resolution hazard does at \num{0.8373} against \num{0.7879}. The admissible
model is honest about its limits: it separates cases at \num{0.8258} AUC and loses only
\num{0.0184} when evaluation moves forward in time rather than across a random split, which is
a smaller drift penalty than an earlier version of this work reported before a date-parsing
defect in its own loader was found and corrected. In the assignment component the proposed
operators improve on their own base algorithm, a reference-point control improves on them
again, and an ablation attributes essentially the whole gain to one of the three proposed
modifications.

Three directions follow. Off-policy evaluation would replace the model-scored counterfactual
used here with an estimator that has a variance bound. No log in this study carries both team
and individual resource identifiers, so whether encoding at individual level is worth its
higher cold-start cost cannot be settled here and needs a log that offers both. And the widest
open question is now sharper than the one this work began with: not whether the cost of
admissibility generalises, since it plainly does not, but what property of a log predicts it.
The five measured here suggest the answer has more to do with what each organisation happened
to record after a decision than with the domain it operates in.


% Shared block. Included by BOTH paper.tex and paper_ieee.tex.
% Do NOT duplicate this text into either driver. The two drivers previously carried
% their own copies of the abstract, conclusion and declarations, and they diverged:
% the IEEE build shipped a code-availability URL that the single-column build did not,
% and both PDFs went out that way on 2026-09-07.
\section*{Data availability}

BPI Challenge 2014 is publicly available from 4TU.ResearchData under DOI
\href{https://doi.org/10.4121/uuid:c3e5d162-0cfd-4bb0-bd82-af5268819c35}%
{10.4121/\allowbreak uuid:\allowbreak c3e5d162-0cfd-\allowbreak 4bb0-bd82-\allowbreak af5268819c35}. The incident-detail file used here
carries DOI \href{https://doi.org/10.4121/uuid:3cfa2260-f5c5-44be-afe1-b70d35288d6d}%
{10.4121/\allowbreak uuid:\allowbreak 3cfa2260-f5c5-\allowbreak 44be-afe1-\allowbreak b70d35288d6d}. No data were collected by the authors.

\section*{Code availability}

All analysis code, including every script required to regenerate every number, table and figure
in this paper, is publicly available at
\href{https://github.com/ashwanth2007/bpi2014-admissibility}{github.com/ashwanth2007/\allowbreak bpi2014-admissibility}. Every experiment is seeded and reproducible: two
invocations with identical arguments produce bitwise-identical metrics.

\section*{Funding}

This research received no specific grant from any funding agency in the public, commercial,
or not-for-profit sectors.

\section*{Author contributions}

% TO BE CONFIRMED WITH THE TEAM BEFORE SUBMISSION. Placeholder, not a claim.

\section*{Competing interests}

The authors declare no competing interests.

\section*{Ethics}

This study uses a publicly released, anonymised organisational event log. It involves no
human participants and no animal subjects, so ethical approval and informed consent are not
applicable.


\bibliography{refs}

\end{document}
